\documentclass[../main.tex]{subfiles}
\begin{document}

\chapter{Multi GPU}
\lessoninfo{
  video={M9, lessons 1--2 (Multi-GPU Hardware; Concurrency on GPUs)},
  sources={Arunkumar et al.\ \cite{arunkumar2017mcmgpu} (required); NVIDIA NVLink page \cite{nvidia2026nvlink}, GPUDirect RDMA documentation \cite{nvidia2026gpudirectrdma}, MPS overview \cite{nvidia2026mps}, Yu et al.\ \cite{yu2022multitenant}, Ye et al.\ \cite{ye2024dlsched} (optional); NVIDIA A100 \cite{nvidia2020ampere} and H100 \cite{nvidia2022hopper} whitepapers},
  textbook={A18 \cite{aamodt2018gpu} ch.~5; PMPP \cite{hwu2022pmpp} ch.~20; CPG \cite{nvidia2024cudaguide} §3.2.8 Asynchronous Concurrent Execution, §3.2.9 Multi-Device System},
  keywords={multi-GPU, monolithic GPU, multi-chip module, GPU module, NUMA, near and far memory, first-touch page placement, page migration, thread block scheduling, NVLink, NVSwitch, RDMA, GPUDirect RDMA, utilization, multi-tenancy, MPS, MIG, concurrent kernel execution, CUDA Graph},
}

The lessons so far have treated one GPU: its programming model
(Lessons 3 and~5), its SMs and memory system (Lessons 4 and~6), and its
virtual memory and warp scheduling (Lesson~7).  Many workloads,
datacenter training of large models above all, need more computing power
than one chip provides, and this lesson asks how GPUs are combined.  The
first part follows the hardware from one die to several dies in a
package, several GPUs in a node and several nodes on a network: the
memory becomes non-uniform, so where a thread block runs and where its
pages live start to matter, and new interconnects (NVLink, NVSwitch,
RDMA) carry the traffic that PCIe cannot.  The second part asks the
opposite question: when one GPU is more than a job needs, how can
several jobs share it?  Streams, the Multi-Process Service and the
Multi-Instance GPU are three answers that differ in how much of the GPU
they partition and how much isolation they give.  The distributed-memory
programming of Lesson~2 returns at the largest scale.

\section{Multi-GPU hardware}
\label{sec:l09-hardware}

The demand for computing power has grown faster than a single GPU,
and GPUs have answered with more chips.\source{M9, multi-GPU
hardware; A18 ch.~5; \cite{arunkumar2017mcmgpu}.}  This section follows
the steps of that scaling, the memory organisation they produce, the
two placement decisions it raises, and the interconnects that join GPUs
inside a node and across nodes.

\subsection{Scaling beyond one chip}

The first way to a faster GPU is more cores on one die.

\begin{definition}[Monolithic GPU, multi-GPU]
  A \term{monolithic GPU}[モノリシックGPU] is a GPU built as a single
  die: all SMs, caches and memory controllers on one piece of
  silicon.\margin{SMs per data-centre GPU: 56 (P100), 80 (V100), 108
  (A100), 132 (H100 SXM5, preliminary); transistors from 15.3 to 80
  billion over the same four generations (AWP Table~4
  \cite{nvidia2020ampere}; HWP Table~3 \cite{nvidia2022hopper}).}
  A \term{multi-GPU}[マルチGPU] system combines several GPU chips, in one
  package, on one board or across boards, to act on one workload.
\end{definition}

Every GPU generation of Lesson~1 added SMs to a monolithic die, and
\cref{fig:l09-scaling} shows where that path ends.  A larger die costs
more, and its \emph{yield}, the fraction of manufactured dies that
work, falls, because every die is more likely to contain a fatal
defect; and the lithography sets a hard maximum die size.\margin{Data-centre
dies are already large: \SI{815}{\milli\metre\squared} (V100),
\SI{826}{\milli\metre\squared} (A100) and \SI{814}{\milli\metre\squared}
(H100) (AWP Table~4 \cite{nvidia2020ampere}; HWP Table~3
\cite{nvidia2022hopper}).  Arunkumar et al.\ assume that a GPU of more than 128 SMs cannot be
built as one die \cite{arunkumar2017mcmgpu}.}  The next step therefore keeps
the dies small and connects several of them, either as separate GPU
chips in several sockets or as several dies on a silicon interposer
(Lesson~4) inside one package, joined by a high-bandwidth switch.  The
last step connects whole GPU boards through PCIe.

\begin{figure}[tb]
  \centering
  % How GPUs scale up (Lesson 9): from one chip to multi-GPU boards.
% Four stages left to right: a small GPU (cores, cache, memory); a
% monolithic GPU with four times the cores; four GPU chips with their own
% memories joined by a high-bandwidth switch inside one package or
% system; several GPU boards joined by PCIe.
% Redraws outline figure m9l1_multigpu.
% Usage: % How GPUs scale up (Lesson 9): from one chip to multi-GPU boards.
% Four stages left to right: a small GPU (cores, cache, memory); a
% monolithic GPU with four times the cores; four GPU chips with their own
% memories joined by a high-bandwidth switch inside one package or
% system; several GPU boards joined by PCIe.
% Redraws outline figure m9l1_multigpu.
% Usage: % How GPUs scale up (Lesson 9): from one chip to multi-GPU boards.
% Four stages left to right: a small GPU (cores, cache, memory); a
% monolithic GPU with four times the cores; four GPU chips with their own
% memories joined by a high-bandwidth switch inside one package or
% system; several GPU boards joined by PCIe.
% Redraws outline figure m9l1_multigpu.
% Usage: \input{figures/l09-scaling}
% Width: body (120 mm); the drawing is about 117 mm wide.
\begin{tikzpicture}[x=1mm, y=1mm, font=\scriptsize\sffamily,
  core/.style={draw, fill=ln-accent!30, minimum width=3mm, minimum height=2.4mm, inner sep=0pt},
  bar/.style={draw, inner sep=0pt, align=center},
  arr/.style={-{Stealth[length=2mm]}, line width=0.8pt, ln-tip},
  stcap/.style={align=center, text width=27mm, anchor=north}]
  % --- stage 1: small GPU (8 cores)
  \foreach \i in {0,...,3} {
    \node[core] at (1.5+\i*3.6, 23) {};
    \node[core] at (1.5+\i*3.6, 16) {};
  }
  \node[bar, fill=ln-box, minimum width=14mm, minimum height=3mm] at (6.9,19.5) {\iflangja{キャッシュ}{cache}};
  \node[bar, fill=ln-muted!20, minimum width=14mm, minimum height=3.5mm] at (6.9,9) {\iflangja{メモリ}{Mem}};
  \node[stcap] at (6.9,5) {\iflangja{小規模 GPU}{small GPU}};
  \draw[arr] (16,17) -- (21,17);
  % --- stage 2: monolithic GPU (16 cores)
  \foreach \i in {0,...,3} {
    \foreach \r in {0,1} {
      \node[core] at (23.5+\i*3.6, 26.5-\r*3.2) {};
      \node[core] at (23.5+\i*3.6, 15.5-\r*3.2) {};
    }
  }
  \node[bar, fill=ln-box, minimum width=14mm, minimum height=3mm] at (28.9,19.5) {\iflangja{キャッシュ}{cache}};
  \node[bar, fill=ln-muted!20, minimum width=14mm, minimum height=3.5mm] at (28.9,7) {\iflangja{メモリ}{Mem}};
  \node[stcap] at (28.9,4) {\iflangja{モノリシック GPU}{monolithic GPU}};
  \draw[arr] (38,17) -- (43,17);
  % --- stage 3: four GPU chips and a switch
  \foreach \x in {0,1} {
    \node[bar, fill=ln-muted!20, minimum width=11mm, minimum height=3mm] at (51+\x*14,29) {\iflangja{メモリ}{Mem}};
    \node[bar, fill=ln-accent!30, minimum width=11mm, minimum height=5mm] at (51+\x*14,23.5) {GPU};
    \node[bar, fill=ln-accent!30, minimum width=11mm, minimum height=5mm] at (51+\x*14,12) {GPU};
    \node[bar, fill=ln-muted!20, minimum width=11mm, minimum height=3mm] at (51+\x*14,6.5) {\iflangja{メモリ}{Mem}};
  }
  \node[bar, fill=ln-caution!25, minimum width=25mm, minimum height=4mm, text width=25mm] at (58,17.75) {\iflangja{高帯域スイッチ}{high-bandwidth switch}};
  \node[stcap] at (58,4) {\iflangja{マルチソケット／マルチチップ GPU}{multi-socket or multi-chip GPUs}};
  \draw[arr] (73,17) -- (78,17);
  % --- stage 4: boards on PCIe
  \node[bar, fill=ln-box, minimum width=22mm, minimum height=6mm] at (91,25) {\iflangja{GPU ボード}{GPU board}};
  \node[bar, fill=ln-caution!25, minimum width=22mm, minimum height=4mm] at (91,17.75) {PCIe};
  \node[bar, fill=ln-box, minimum width=22mm, minimum height=6mm] at (91,10.5) {\iflangja{GPU ボード}{GPU board}};
  \node[stcap] at (91,4) {\iflangja{マルチ GPU ボード}{multi-GPU boards}};
\end{tikzpicture}

% Width: body (120 mm); the drawing is about 117 mm wide.
\begin{tikzpicture}[x=1mm, y=1mm, font=\scriptsize\sffamily,
  core/.style={draw, fill=ln-accent!30, minimum width=3mm, minimum height=2.4mm, inner sep=0pt},
  bar/.style={draw, inner sep=0pt, align=center},
  arr/.style={-{Stealth[length=2mm]}, line width=0.8pt, ln-tip},
  stcap/.style={align=center, text width=27mm, anchor=north}]
  % --- stage 1: small GPU (8 cores)
  \foreach \i in {0,...,3} {
    \node[core] at (1.5+\i*3.6, 23) {};
    \node[core] at (1.5+\i*3.6, 16) {};
  }
  \node[bar, fill=ln-box, minimum width=14mm, minimum height=3mm] at (6.9,19.5) {\iflangja{キャッシュ}{cache}};
  \node[bar, fill=ln-muted!20, minimum width=14mm, minimum height=3.5mm] at (6.9,9) {\iflangja{メモリ}{Mem}};
  \node[stcap] at (6.9,5) {\iflangja{小規模 GPU}{small GPU}};
  \draw[arr] (16,17) -- (21,17);
  % --- stage 2: monolithic GPU (16 cores)
  \foreach \i in {0,...,3} {
    \foreach \r in {0,1} {
      \node[core] at (23.5+\i*3.6, 26.5-\r*3.2) {};
      \node[core] at (23.5+\i*3.6, 15.5-\r*3.2) {};
    }
  }
  \node[bar, fill=ln-box, minimum width=14mm, minimum height=3mm] at (28.9,19.5) {\iflangja{キャッシュ}{cache}};
  \node[bar, fill=ln-muted!20, minimum width=14mm, minimum height=3.5mm] at (28.9,7) {\iflangja{メモリ}{Mem}};
  \node[stcap] at (28.9,4) {\iflangja{モノリシック GPU}{monolithic GPU}};
  \draw[arr] (38,17) -- (43,17);
  % --- stage 3: four GPU chips and a switch
  \foreach \x in {0,1} {
    \node[bar, fill=ln-muted!20, minimum width=11mm, minimum height=3mm] at (51+\x*14,29) {\iflangja{メモリ}{Mem}};
    \node[bar, fill=ln-accent!30, minimum width=11mm, minimum height=5mm] at (51+\x*14,23.5) {GPU};
    \node[bar, fill=ln-accent!30, minimum width=11mm, minimum height=5mm] at (51+\x*14,12) {GPU};
    \node[bar, fill=ln-muted!20, minimum width=11mm, minimum height=3mm] at (51+\x*14,6.5) {\iflangja{メモリ}{Mem}};
  }
  \node[bar, fill=ln-caution!25, minimum width=25mm, minimum height=4mm, text width=25mm] at (58,17.75) {\iflangja{高帯域スイッチ}{high-bandwidth switch}};
  \node[stcap] at (58,4) {\iflangja{マルチソケット／マルチチップ GPU}{multi-socket or multi-chip GPUs}};
  \draw[arr] (73,17) -- (78,17);
  % --- stage 4: boards on PCIe
  \node[bar, fill=ln-box, minimum width=22mm, minimum height=6mm] at (91,25) {\iflangja{GPU ボード}{GPU board}};
  \node[bar, fill=ln-caution!25, minimum width=22mm, minimum height=4mm] at (91,17.75) {PCIe};
  \node[bar, fill=ln-box, minimum width=22mm, minimum height=6mm] at (91,10.5) {\iflangja{GPU ボード}{GPU board}};
  \node[stcap] at (91,4) {\iflangja{マルチ GPU ボード}{multi-GPU boards}};
\end{tikzpicture}

% Width: body (120 mm); the drawing is about 117 mm wide.
\begin{tikzpicture}[x=1mm, y=1mm, font=\scriptsize\sffamily,
  core/.style={draw, fill=ln-accent!30, minimum width=3mm, minimum height=2.4mm, inner sep=0pt},
  bar/.style={draw, inner sep=0pt, align=center},
  arr/.style={-{Stealth[length=2mm]}, line width=0.8pt, ln-tip},
  stcap/.style={align=center, text width=27mm, anchor=north}]
  % --- stage 1: small GPU (8 cores)
  \foreach \i in {0,...,3} {
    \node[core] at (1.5+\i*3.6, 23) {};
    \node[core] at (1.5+\i*3.6, 16) {};
  }
  \node[bar, fill=ln-box, minimum width=14mm, minimum height=3mm] at (6.9,19.5) {\iflangja{キャッシュ}{cache}};
  \node[bar, fill=ln-muted!20, minimum width=14mm, minimum height=3.5mm] at (6.9,9) {\iflangja{メモリ}{Mem}};
  \node[stcap] at (6.9,5) {\iflangja{小規模 GPU}{small GPU}};
  \draw[arr] (16,17) -- (21,17);
  % --- stage 2: monolithic GPU (16 cores)
  \foreach \i in {0,...,3} {
    \foreach \r in {0,1} {
      \node[core] at (23.5+\i*3.6, 26.5-\r*3.2) {};
      \node[core] at (23.5+\i*3.6, 15.5-\r*3.2) {};
    }
  }
  \node[bar, fill=ln-box, minimum width=14mm, minimum height=3mm] at (28.9,19.5) {\iflangja{キャッシュ}{cache}};
  \node[bar, fill=ln-muted!20, minimum width=14mm, minimum height=3.5mm] at (28.9,7) {\iflangja{メモリ}{Mem}};
  \node[stcap] at (28.9,4) {\iflangja{モノリシック GPU}{monolithic GPU}};
  \draw[arr] (38,17) -- (43,17);
  % --- stage 3: four GPU chips and a switch
  \foreach \x in {0,1} {
    \node[bar, fill=ln-muted!20, minimum width=11mm, minimum height=3mm] at (51+\x*14,29) {\iflangja{メモリ}{Mem}};
    \node[bar, fill=ln-accent!30, minimum width=11mm, minimum height=5mm] at (51+\x*14,23.5) {GPU};
    \node[bar, fill=ln-accent!30, minimum width=11mm, minimum height=5mm] at (51+\x*14,12) {GPU};
    \node[bar, fill=ln-muted!20, minimum width=11mm, minimum height=3mm] at (51+\x*14,6.5) {\iflangja{メモリ}{Mem}};
  }
  \node[bar, fill=ln-caution!25, minimum width=25mm, minimum height=4mm, text width=25mm] at (58,17.75) {\iflangja{高帯域スイッチ}{high-bandwidth switch}};
  \node[stcap] at (58,4) {\iflangja{マルチソケット／マルチチップ GPU}{multi-socket or multi-chip GPUs}};
  \draw[arr] (73,17) -- (78,17);
  % --- stage 4: boards on PCIe
  \node[bar, fill=ln-box, minimum width=22mm, minimum height=6mm] at (91,25) {\iflangja{GPU ボード}{GPU board}};
  \node[bar, fill=ln-caution!25, minimum width=22mm, minimum height=4mm] at (91,17.75) {PCIe};
  \node[bar, fill=ln-box, minimum width=22mm, minimum height=6mm] at (91,10.5) {\iflangja{GPU ボード}{GPU board}};
  \node[stcap] at (91,4) {\iflangja{マルチ GPU ボード}{multi-GPU boards}};
\end{tikzpicture}

  \caption{Four steps of GPU scaling: a small GPU; a monolithic GPU with
    more cores on one die; several GPU chips, each with its own memory,
    joined by a high-bandwidth switch (multi-socket or multi-chip); and
    several GPU boards joined by PCIe.}
  \label{fig:l09-scaling}
\end{figure}

\begin{definition}[Multi-chip module, GPU module, chiplet]
  A \term{multi-chip module (MCM)}[マルチチップモジュール]\index{MCM|see{multi-chip module (MCM)}} is a package
  that holds several dies connected by short on-package links.  In the
  MCM-GPU of Arunkumar et al.\ each die is a
  \term{GPU module (GPM)}[GPUモジュール]\index{GPM|see{GPU module (GPM)}}, a smaller GPU with its own SMs,
  L2 slice and DRAM partition, and the package of GPMs appears to
  software as one larger logical GPU \cite{arunkumar2017mcmgpu}.  A die
  designed to be combined with others in this way is called a
  \term{chiplet}[チップレット].
\end{definition}

The paper calls this \enquote{package-level integration of multiple GPU
modules to build larger logical GPUs}, and argues that it allows
\enquote{continuous performance scaling beyond Moore's law}, because
each GPM stays small enough to be manufactured with good yield
\cite{arunkumar2017mcmgpu}.\sidenote{The MCM-GPU studied in the paper
has four GPMs of 64 SMs each (256 SMs), \SI{3}{\tera\byte\per\second}
of DRAM bandwidth in total, \SI{768}{\giga\byte\per\second} per local
partition.  Its optimised version is reported as \enquote{45.5\,\%
faster than the largest implementable monolithic GPU}, within 10\,\% of
an unbuildable monolithic GPU of the same size, and 26.8\,\% faster than
a multi-GPU system with the same SMs and DRAM bandwidth
\cite{arunkumar2017mcmgpu}.}

\subsection{One logical GPU with non-uniform memory}

To the programmer, a package or node of several GPUs joined by
high-speed links (\cref{fig:l09-numa}) is still a large number of SMs,
and all memories are reachable as global memory, one shared address
space.  Each GPU keeps its own DRAM, typically HBM stacks (Lesson~1)
for the bandwidth its SMs demand, while the I/O bandwidth to the rest of
the system is shared by all of them.  What changes is the cost of an
access.

\begin{definition}[NUMA, near and far memory]
  A memory system has
  \term{non-uniform memory access (NUMA)}[不均一メモリアクセス（NUMA）]
  when one address space is spread over several memories and the latency
  and bandwidth of an access depend on which memory holds the address
  relative to the processor that issues it.
  \index{NUMA|see{non-uniform memory access (NUMA)}}%
  The two cases are called
  \term{near memory / far memory}[近接メモリ／遠隔メモリ]: near memory
  is attached to the GPU that issues the access, far memory to another
  GPU, so that the access must cross the inter-GPU link.
\end{definition}

\begin{figure}[tb]
  \centering
  % A four-GPU NUMA system (Lesson 9).  Four GPUs, each with its own
% memory on the outer side, fully connected by high-speed links (the six
% links of a 2x2 all-to-all).  A thread block on GPU 0 reaches a page in
% its own memory directly (near access) and a page in GPU 3's memory over
% a link (far access).
% Redraws outline figures m9l1_nvlink and m9l1_tb_s3 (simplified).
% Usage: % A four-GPU NUMA system (Lesson 9).  Four GPUs, each with its own
% memory on the outer side, fully connected by high-speed links (the six
% links of a 2x2 all-to-all).  A thread block on GPU 0 reaches a page in
% its own memory directly (near access) and a page in GPU 3's memory over
% a link (far access).
% Redraws outline figures m9l1_nvlink and m9l1_tb_s3 (simplified).
% Usage: % A four-GPU NUMA system (Lesson 9).  Four GPUs, each with its own
% memory on the outer side, fully connected by high-speed links (the six
% links of a 2x2 all-to-all).  A thread block on GPU 0 reaches a page in
% its own memory directly (near access) and a page in GPU 3's memory over
% a link (far access).
% Redraws outline figures m9l1_nvlink and m9l1_tb_s3 (simplified).
% Usage: \input{figures/l09-numa}
% Width: body (120 mm); the drawing is about 112 mm wide.
\begin{tikzpicture}[x=1mm, y=1mm, font=\scriptsize\sffamily,
  gpu/.style={draw, fill=ln-accent!30, minimum width=20mm, minimum height=10mm, align=center, inner sep=1pt},
  mem/.style={draw, fill=ln-muted!20, minimum width=7mm, minimum height=10mm, align=center, inner sep=1pt},
  lnk/.style={line width=1.6pt, ln-source!70},
  acc/.style={-{Stealth[length=2.2mm]}, line width=1pt}]
  \node[gpu, label={[anchor=north, inner sep=1pt]north:GPU 0}] (g0) at (20,32) {};
  \node[gpu] (g1) at (60,32) {GPU 1};
  \node[gpu] (g2) at (20,8)  {GPU 2};
  \node[gpu] (g3) at (60,8)  {GPU 3};
  \node[mem] (m0) at (3,32) {\iflangja{メモリ}{Mem}\\0};
  \node[mem] (m1) at (77,32) {\iflangja{メモリ}{Mem}\\1};
  \node[mem] (m2) at (3,8) {\iflangja{メモリ}{Mem}\\2};
  \node[mem] (m3) at (77,8) {\iflangja{メモリ}{Mem}\\3};
  \foreach \g/\m in {g0/m0, g1/m1, g2/m2, g3/m3} \draw[line width=1.2pt] (\g) -- (\m);
  \draw[lnk] (g0) -- (g1);
  \draw[lnk] (g2) -- (g3);
  \draw[lnk] (g0) -- (g2);
  \draw[lnk] (g1) -- (g3);
  \draw[lnk] (g0) -- (g3);
  \draw[lnk] (g1) -- (g2);
  % thread block on GPU 0 and its two accesses
  \node[draw, fill=ln-tip!35, inner sep=1pt] (tb) at (25,29.5) {TB};
  \draw[acc, ln-tip] (tb.west) -- node[above, pos=0.45] {\iflangja{近接}{near}} (m0.east|-tb);
  \draw[acc, ln-caution, dashed] (tb.south east) -- (57,12.5) -- (m3.west|-0,12.5);
  \node[ln-caution, anchor=south west, inner sep=1pt] at (61,13.5) {\iflangja{遠隔アクセス}{far access}};
  % legend
  \draw[lnk] (86,26) -- (92,26);
  \node[anchor=west, align=left, text width=18mm] at (92,26) {\iflangja{高速リンク（例：NVLink）}{high-speed link (e.g.\ NVLink)}};
\end{tikzpicture}

% Width: body (120 mm); the drawing is about 112 mm wide.
\begin{tikzpicture}[x=1mm, y=1mm, font=\scriptsize\sffamily,
  gpu/.style={draw, fill=ln-accent!30, minimum width=20mm, minimum height=10mm, align=center, inner sep=1pt},
  mem/.style={draw, fill=ln-muted!20, minimum width=7mm, minimum height=10mm, align=center, inner sep=1pt},
  lnk/.style={line width=1.6pt, ln-source!70},
  acc/.style={-{Stealth[length=2.2mm]}, line width=1pt}]
  \node[gpu, label={[anchor=north, inner sep=1pt]north:GPU 0}] (g0) at (20,32) {};
  \node[gpu] (g1) at (60,32) {GPU 1};
  \node[gpu] (g2) at (20,8)  {GPU 2};
  \node[gpu] (g3) at (60,8)  {GPU 3};
  \node[mem] (m0) at (3,32) {\iflangja{メモリ}{Mem}\\0};
  \node[mem] (m1) at (77,32) {\iflangja{メモリ}{Mem}\\1};
  \node[mem] (m2) at (3,8) {\iflangja{メモリ}{Mem}\\2};
  \node[mem] (m3) at (77,8) {\iflangja{メモリ}{Mem}\\3};
  \foreach \g/\m in {g0/m0, g1/m1, g2/m2, g3/m3} \draw[line width=1.2pt] (\g) -- (\m);
  \draw[lnk] (g0) -- (g1);
  \draw[lnk] (g2) -- (g3);
  \draw[lnk] (g0) -- (g2);
  \draw[lnk] (g1) -- (g3);
  \draw[lnk] (g0) -- (g3);
  \draw[lnk] (g1) -- (g2);
  % thread block on GPU 0 and its two accesses
  \node[draw, fill=ln-tip!35, inner sep=1pt] (tb) at (25,29.5) {TB};
  \draw[acc, ln-tip] (tb.west) -- node[above, pos=0.45] {\iflangja{近接}{near}} (m0.east|-tb);
  \draw[acc, ln-caution, dashed] (tb.south east) -- (57,12.5) -- (m3.west|-0,12.5);
  \node[ln-caution, anchor=south west, inner sep=1pt] at (61,13.5) {\iflangja{遠隔アクセス}{far access}};
  % legend
  \draw[lnk] (86,26) -- (92,26);
  \node[anchor=west, align=left, text width=18mm] at (92,26) {\iflangja{高速リンク（例：NVLink）}{high-speed link (e.g.\ NVLink)}};
\end{tikzpicture}

% Width: body (120 mm); the drawing is about 112 mm wide.
\begin{tikzpicture}[x=1mm, y=1mm, font=\scriptsize\sffamily,
  gpu/.style={draw, fill=ln-accent!30, minimum width=20mm, minimum height=10mm, align=center, inner sep=1pt},
  mem/.style={draw, fill=ln-muted!20, minimum width=7mm, minimum height=10mm, align=center, inner sep=1pt},
  lnk/.style={line width=1.6pt, ln-source!70},
  acc/.style={-{Stealth[length=2.2mm]}, line width=1pt}]
  \node[gpu, label={[anchor=north, inner sep=1pt]north:GPU 0}] (g0) at (20,32) {};
  \node[gpu] (g1) at (60,32) {GPU 1};
  \node[gpu] (g2) at (20,8)  {GPU 2};
  \node[gpu] (g3) at (60,8)  {GPU 3};
  \node[mem] (m0) at (3,32) {\iflangja{メモリ}{Mem}\\0};
  \node[mem] (m1) at (77,32) {\iflangja{メモリ}{Mem}\\1};
  \node[mem] (m2) at (3,8) {\iflangja{メモリ}{Mem}\\2};
  \node[mem] (m3) at (77,8) {\iflangja{メモリ}{Mem}\\3};
  \foreach \g/\m in {g0/m0, g1/m1, g2/m2, g3/m3} \draw[line width=1.2pt] (\g) -- (\m);
  \draw[lnk] (g0) -- (g1);
  \draw[lnk] (g2) -- (g3);
  \draw[lnk] (g0) -- (g2);
  \draw[lnk] (g1) -- (g3);
  \draw[lnk] (g0) -- (g3);
  \draw[lnk] (g1) -- (g2);
  % thread block on GPU 0 and its two accesses
  \node[draw, fill=ln-tip!35, inner sep=1pt] (tb) at (25,29.5) {TB};
  \draw[acc, ln-tip] (tb.west) -- node[above, pos=0.45] {\iflangja{近接}{near}} (m0.east|-tb);
  \draw[acc, ln-caution, dashed] (tb.south east) -- (57,12.5) -- (m3.west|-0,12.5);
  \node[ln-caution, anchor=south west, inner sep=1pt] at (61,13.5) {\iflangja{遠隔アクセス}{far access}};
  % legend
  \draw[lnk] (86,26) -- (92,26);
  \node[anchor=west, align=left, text width=18mm] at (92,26) {\iflangja{高速リンク（例：NVLink）}{high-speed link (e.g.\ NVLink)}};
\end{tikzpicture}

  \caption{A four-GPU NUMA system.  Each GPU has its own memory, and the
    six links connect every pair.  A thread block on GPU~0 reaches its
    own memory directly (near access) and the memory of GPU~3 over a link
    (far access), at lower bandwidth and higher latency.}
  \label{fig:l09-numa}
\end{figure}

\begin{example}[Near and far, in numbers]
  \label{ex:l09-near-far}
  A kernel on an H100 reads a \SI{1}{\gibi\byte} buffer
  (\SI{1.074}{\giga\byte}) once.  From its own HBM3 at the preliminary
  \SI{3}{\tera\byte\per\second} the transfer takes at least
  \SI{0.36}{\milli\second}.  If the buffer lives on a peer GPU, it
  arrives over NVLink~4 at \SI{450}{\giga\byte\per\second} per direction,
  half of the \SI{900}{\giga\byte\per\second} total, in at least
  \SI{2.4}{\milli\second}; over PCIe Gen~5 at
  \SI{64}{\giga\byte\per\second} per direction it takes at least
  \SI{16.8}{\milli\second} \cite{nvidia2022hopper}.  The far access is
  $6.7\times$ slower than the near one over NVLink and $47\times$ slower
  over PCIe, before any difference in latency is counted.
\end{example}

\margin{The figures in \cref{ex:l09-near-far} are peak bandwidths; a
real transfer reaches less, and a pattern of small far accesses suffers
the link latency as well.}%
The ratio, not the absolute figure, is what NUMA means for a program:
data should be near the SMs that use it.  Two questions follow, and in
this chapter they are asked for thread blocks and pages.

\subsection{Where to run a block and where to put a page}

\begin{definition}[Thread block scheduling, first-touch page placement]
  The decision on which GPU (or GPM), and on which of its SMs, each
  thread block of a kernel runs is called
  \term{thread block scheduling}[スレッドブロックスケジューリング].  Under
  \term{first-touch page placement}[ファーストタッチ配置] a page is
  allocated in the memory near the GPU whose thread block accesses it
  first.
\end{definition}

First-touch needs no knowledge of the program: the page goes where it is
first wanted.  It works well as long as each page is used by blocks on
the same GPU; it fails when a later block that runs elsewhere needs the
same page.  Suppose blocks are dealt out round-robin, one per GPU, and
each brings its pages into its GPU's memory by first touch.  A block
scheduled later on GPU~0 that shares a page with the block on GPU~3
finds that page already placed in GPU~3's memory, and every access to it
is a far access (\cref{fig:l09-numa}).  There are two remedies: schedule
the block on the GPU where its data lives, or move the data.

\begin{definition}[Page migration]
  The move of a page from one GPU's memory to another's while the
  program runs is called \term{page migration}[ページマイグレーション]; it is typically done when the
  page has received enough far accesses from that other GPU to repay
  the cost of the copy.
\end{definition}

\margin{Migration is the multi-GPU form of the page movement of
Lesson~7, where unified memory moves pages between host and device on a
page fault; Li et al.\ study the address translation of such
multi-GPU systems \cite{li2021tlb}.}%
A migration is a copy of a whole page over the link plus an update of the
page tables, so a policy that migrates on the first far access can do
more harm than good when two GPUs take turns at the same page.

\begin{example}[Round-robin versus grouped block scheduling]
  \label{ex:l09-placement}
  A kernel of eight thread blocks TB0--TB7 runs on four GPUs, two blocks
  per GPU in two waves.  Block TB$i$ reads page P$i$ and then page
  P$(i{+}1)$, as in a one-dimensional stencil in which neighbouring
  blocks share a boundary page; pages are placed by first touch, and
  within a wave the blocks touch their pages in index order.
  \Cref{tab:l09-placement} traces two schedules.  Round-robin puts TB$i$
  on GPU $i \bmod 4$, so neighbours always sit on different GPUs: 7 of
  the 16 page accesses are far.  Grouping contiguous blocks, TB0 and TB1
  on GPU~0, TB2 and TB3 on GPU~1 and so on, leaves a shared page far only
  at the boundary between two groups: 3 of 16.
\end{example}

\begin{table}[tb]
  \centering
  \small
  \setlength{\tabcolsep}{4pt}
  \caption{Trace of \cref{ex:l09-placement}.  Each entry is a block, its
    GPU, and its two page accesses; \enquote{far} marks a page that an
    earlier block has already placed on another GPU.}
  \label{tab:l09-placement}
  \begin{tabular}{@{}l l l l l l@{}}
    \toprule
    & \multicolumn{2}{l}{Round-robin} & & \multicolumn{2}{l}{Grouped} \\
    \cmidrule{2-3}\cmidrule{5-6}
    Wave & Block, GPU & Pages & & Block, GPU & Pages \\
    \midrule
    1 & TB0, G0 & P0 new, P1 new & & TB0, G0 & P0 new, P1 new \\
      & TB1, G1 & P1 \textbf{far}, P2 new & & TB2, G1 & P2 new, P3 new \\
      & TB2, G2 & P2 \textbf{far}, P3 new & & TB4, G2 & P4 new, P5 new \\
      & TB3, G3 & P3 \textbf{far}, P4 new & & TB6, G3 & P6 new, P7 new \\
    \addlinespace
    2 & TB4, G0 & P4 \textbf{far}, P5 new & & TB1, G0 & P1 near, P2 \textbf{far} \\
      & TB5, G1 & P5 \textbf{far}, P6 new & & TB3, G1 & P3 near, P4 \textbf{far} \\
      & TB6, G2 & P6 \textbf{far}, P7 new & & TB5, G2 & P5 near, P6 \textbf{far} \\
      & TB7, G3 & P7 \textbf{far}, P8 new & & TB7, G3 & P7 near, P8 new \\
    \midrule
    \multicolumn{3}{@{}l}{Far accesses: 7 of 16} & & \multicolumn{2}{l}{3 of 16} \\
    \bottomrule
  \end{tabular}
\end{table}

\caution{Lecture vs paper: this chapter, like the lecture, reduces the
problem to first-touch placement and locality-aware block scheduling.
The MCM-GPU paper proposes three optimisations together: a GPM-side
\enquote{L1.5} cache that holds only remote data, \enquote{distributed
and batched} CTA scheduling that sends a group of contiguous CTAs to
one GPM, and first-touch page allocation; combined they give a 22.8\,\%
speed-up and a $5\times$ reduction of inter-GPM traffic over the basic
MCM-GPU \cite{arunkumar2017mcmgpu}.}%
The grouped schedule of \cref{ex:l09-placement} is the idea of the
paper's distributed scheduler: contiguous blocks tend to share data, so
they should share a GPU.  It also pays across kernel launches: a kernel
launched again in a loop sends the same block indices to the same GPUs,
whose memories already hold the pages that first touch put there
\cite{arunkumar2017mcmgpu}.

\subsection{NVLink and NVSwitch}

The link between the GPUs decides how much a far access costs.

\begin{definition}[Interconnect, NVLink, NVSwitch]
  An \term{interconnect}[インターコネクト] is the network of links (and
  switches) over which processors and memories exchange data.
  NVIDIA's \term{NVLink}[NVLink] is a point-to-point GPU-to-GPU
  interconnect and protocol, an alternative to PCIe for traffic between
  GPUs.\margin{NVLink per GPU, by generation: \SI{900}{\giga\byte\per\second}
  and 18 links (4th, Hopper); \SI{1800}{\giga\byte\per\second} and 18
  links (5th, Blackwell); \SI{3000}{\giga\byte\per\second} and 36 links
  (6th, Vera Rubin, preliminary).  An NVLink Switch domain spans 8 GPUs,
  or 72 in an NVL72 rack \cite{nvidia2026nvlink}.}  An \term{NVSwitch}[NVSwitch] is a switch chip with many NVLink
  ports; with NVSwitches between them, every GPU of a node reaches every
  other at full link bandwidth.
\end{definition}


\begin{figure}[tb]
  \centering
  % GPU-to-GPU connections (Lesson 9).  (a) A CPU and two GPUs that talk
% only through a PCIe switch; (b) the same with a direct NVLink between
% the GPUs; (c) eight GPUs, each linked to every NVSwitch chip, so that
% any pair of GPUs is two link hops apart at full link bandwidth.
% Redraws outline figure m9l1_nvlink1.
% Usage: % GPU-to-GPU connections (Lesson 9).  (a) A CPU and two GPUs that talk
% only through a PCIe switch; (b) the same with a direct NVLink between
% the GPUs; (c) eight GPUs, each linked to every NVSwitch chip, so that
% any pair of GPUs is two link hops apart at full link bandwidth.
% Redraws outline figure m9l1_nvlink1.
% Usage: % GPU-to-GPU connections (Lesson 9).  (a) A CPU and two GPUs that talk
% only through a PCIe switch; (b) the same with a direct NVLink between
% the GPUs; (c) eight GPUs, each linked to every NVSwitch chip, so that
% any pair of GPUs is two link hops apart at full link bandwidth.
% Redraws outline figure m9l1_nvlink1.
% Usage: \input{figures/l09-interconnect}
% Width: body (120 mm); the drawing is about 118 mm wide.
\begin{tikzpicture}[x=1mm, y=1mm, font=\scriptsize\sffamily,
  gpu/.style={draw, fill=ln-accent!30, minimum width=9mm, minimum height=5mm, inner sep=1pt},
  cpu/.style={draw, fill=ln-box, minimum width=9mm, minimum height=5mm, inner sep=1pt},
  sw/.style={draw, fill=ln-muted!60, text=white, minimum height=4.5mm, inner sep=1pt, align=center},
  pcie/.style={{Stealth[length=1.6mm]}-{Stealth[length=1.6mm]}, line width=0.6pt, ln-muted},
  nvl/.style={{Stealth[length=1.6mm]}-{Stealth[length=1.6mm]}, line width=1.4pt, ln-source!70}]
  % ---- (a) PCIe only
  \node[cpu] (c) at (14,30) {CPU};
  \node[sw, minimum width=22mm] (s) at (14,20) {\iflangja{PCIe スイッチ}{PCIe switch}};
  \node[gpu] (ga) at (6,8) {GPU};
  \node[gpu] (gb) at (22,8) {GPU};
  \draw[pcie] (c) -- (s); \draw[pcie] (ga) -- (ga|-s.south); \draw[pcie] (gb) -- (gb|-s.south);
  \node at (14,0) {(a)};
  % ---- (b) PCIe + NVLink
  \node[cpu] (c2) at (46,30) {CPU};
  \node[sw, minimum width=22mm] (s2) at (46,20) {\iflangja{PCIe スイッチ}{PCIe switch}};
  \node[gpu] (gc) at (38,8) {GPU};
  \node[gpu] (gd) at (54,8) {GPU};
  \draw[pcie] (c2) -- (s2); \draw[pcie] (gc) -- (gc|-s2.south); \draw[pcie] (gd) -- (gd|-s2.south);
  \draw[nvl] (gc) -- (gd);
  \node at (46,0) {(b)};
  % ---- (c) NVSwitch: links first, switches drawn on top of them
  \foreach \k in {0,...,3} {
    \node[gpu] (L\k) at (77,4+\k*8.6) {GPU};
    \node[gpu] (R\k) at (115,4+\k*8.6) {GPU};
    \draw[ln-source!70, line width=0.8pt] (L\k.east) -- (R\k.west);
  }
  \node[sw, minimum width=4.5mm, minimum height=30mm] (n1) at (91,17) {};
  \node[sw, minimum width=4.5mm, minimum height=30mm] (n2) at (101,17) {};
  \node[text=ln-muted] at (96,35) {NVSwitch};
  \node at (96,0) {(c)};
  % legend
  \draw[pcie] (6,-5) -- (12,-5); \node[anchor=west] at (12,-5) {PCIe};
  \draw[nvl] (30,-5) -- (36,-5); \node[anchor=west] at (36,-5) {NVLink};
\end{tikzpicture}

% Width: body (120 mm); the drawing is about 118 mm wide.
\begin{tikzpicture}[x=1mm, y=1mm, font=\scriptsize\sffamily,
  gpu/.style={draw, fill=ln-accent!30, minimum width=9mm, minimum height=5mm, inner sep=1pt},
  cpu/.style={draw, fill=ln-box, minimum width=9mm, minimum height=5mm, inner sep=1pt},
  sw/.style={draw, fill=ln-muted!60, text=white, minimum height=4.5mm, inner sep=1pt, align=center},
  pcie/.style={{Stealth[length=1.6mm]}-{Stealth[length=1.6mm]}, line width=0.6pt, ln-muted},
  nvl/.style={{Stealth[length=1.6mm]}-{Stealth[length=1.6mm]}, line width=1.4pt, ln-source!70}]
  % ---- (a) PCIe only
  \node[cpu] (c) at (14,30) {CPU};
  \node[sw, minimum width=22mm] (s) at (14,20) {\iflangja{PCIe スイッチ}{PCIe switch}};
  \node[gpu] (ga) at (6,8) {GPU};
  \node[gpu] (gb) at (22,8) {GPU};
  \draw[pcie] (c) -- (s); \draw[pcie] (ga) -- (ga|-s.south); \draw[pcie] (gb) -- (gb|-s.south);
  \node at (14,0) {(a)};
  % ---- (b) PCIe + NVLink
  \node[cpu] (c2) at (46,30) {CPU};
  \node[sw, minimum width=22mm] (s2) at (46,20) {\iflangja{PCIe スイッチ}{PCIe switch}};
  \node[gpu] (gc) at (38,8) {GPU};
  \node[gpu] (gd) at (54,8) {GPU};
  \draw[pcie] (c2) -- (s2); \draw[pcie] (gc) -- (gc|-s2.south); \draw[pcie] (gd) -- (gd|-s2.south);
  \draw[nvl] (gc) -- (gd);
  \node at (46,0) {(b)};
  % ---- (c) NVSwitch: links first, switches drawn on top of them
  \foreach \k in {0,...,3} {
    \node[gpu] (L\k) at (77,4+\k*8.6) {GPU};
    \node[gpu] (R\k) at (115,4+\k*8.6) {GPU};
    \draw[ln-source!70, line width=0.8pt] (L\k.east) -- (R\k.west);
  }
  \node[sw, minimum width=4.5mm, minimum height=30mm] (n1) at (91,17) {};
  \node[sw, minimum width=4.5mm, minimum height=30mm] (n2) at (101,17) {};
  \node[text=ln-muted] at (96,35) {NVSwitch};
  \node at (96,0) {(c)};
  % legend
  \draw[pcie] (6,-5) -- (12,-5); \node[anchor=west] at (12,-5) {PCIe};
  \draw[nvl] (30,-5) -- (36,-5); \node[anchor=west] at (36,-5) {NVLink};
\end{tikzpicture}

% Width: body (120 mm); the drawing is about 118 mm wide.
\begin{tikzpicture}[x=1mm, y=1mm, font=\scriptsize\sffamily,
  gpu/.style={draw, fill=ln-accent!30, minimum width=9mm, minimum height=5mm, inner sep=1pt},
  cpu/.style={draw, fill=ln-box, minimum width=9mm, minimum height=5mm, inner sep=1pt},
  sw/.style={draw, fill=ln-muted!60, text=white, minimum height=4.5mm, inner sep=1pt, align=center},
  pcie/.style={{Stealth[length=1.6mm]}-{Stealth[length=1.6mm]}, line width=0.6pt, ln-muted},
  nvl/.style={{Stealth[length=1.6mm]}-{Stealth[length=1.6mm]}, line width=1.4pt, ln-source!70}]
  % ---- (a) PCIe only
  \node[cpu] (c) at (14,30) {CPU};
  \node[sw, minimum width=22mm] (s) at (14,20) {\iflangja{PCIe スイッチ}{PCIe switch}};
  \node[gpu] (ga) at (6,8) {GPU};
  \node[gpu] (gb) at (22,8) {GPU};
  \draw[pcie] (c) -- (s); \draw[pcie] (ga) -- (ga|-s.south); \draw[pcie] (gb) -- (gb|-s.south);
  \node at (14,0) {(a)};
  % ---- (b) PCIe + NVLink
  \node[cpu] (c2) at (46,30) {CPU};
  \node[sw, minimum width=22mm] (s2) at (46,20) {\iflangja{PCIe スイッチ}{PCIe switch}};
  \node[gpu] (gc) at (38,8) {GPU};
  \node[gpu] (gd) at (54,8) {GPU};
  \draw[pcie] (c2) -- (s2); \draw[pcie] (gc) -- (gc|-s2.south); \draw[pcie] (gd) -- (gd|-s2.south);
  \draw[nvl] (gc) -- (gd);
  \node at (46,0) {(b)};
  % ---- (c) NVSwitch: links first, switches drawn on top of them
  \foreach \k in {0,...,3} {
    \node[gpu] (L\k) at (77,4+\k*8.6) {GPU};
    \node[gpu] (R\k) at (115,4+\k*8.6) {GPU};
    \draw[ln-source!70, line width=0.8pt] (L\k.east) -- (R\k.west);
  }
  \node[sw, minimum width=4.5mm, minimum height=30mm] (n1) at (91,17) {};
  \node[sw, minimum width=4.5mm, minimum height=30mm] (n2) at (101,17) {};
  \node[text=ln-muted] at (96,35) {NVSwitch};
  \node at (96,0) {(c)};
  % legend
  \draw[pcie] (6,-5) -- (12,-5); \node[anchor=west] at (12,-5) {PCIe};
  \draw[nvl] (30,-5) -- (36,-5); \node[anchor=west] at (36,-5) {NVLink};
\end{tikzpicture}

  \caption{(a) Two GPUs that communicate only through the PCIe switch
    they share with the CPU.  (b) The same GPUs with a direct NVLink.
    (c) Eight GPUs, each linked to every NVSwitch chip, so that every
    pair of GPUs is connected through a switch.}
  \label{fig:l09-interconnect}
\end{figure}

Without NVLink, data between two GPUs crosses the PCIe switch that they
share with the CPU (\cref{fig:l09-interconnect}a); a direct NVLink
bypasses it (b).  Fourth-generation NVLink, in the H100, provides
\SI{900}{\giga\byte\per\second} per GPU over 18 links, which the
whitepaper quotes as $7\times$ the bandwidth of PCIe Gen~5
\cite{nvidia2022hopper}; against the H100's \SI{3}{\tera\byte\per\second}
of HBM3 bandwidth, NVLink offers 30\,\% of the local memory bandwidth
(both directions counted).\caution{Lecture vs reality: the lecture calls
this about a quarter.  Per direction NVLink~4 carries
\SI{450}{\giga\byte\per\second} (18 links of
\SI{25}{\giga\byte\per\second}), 15\,\% of HBM \cite{nvidia2022hopper}.}  Point-to-point links do not scale, however: a
GPU has a fixed number of links, and connecting every pair of $n$ GPUs
directly needs $n(n-1)/2$ links.  NVSwitch chips solve this
(\cref{fig:l09-interconnect}c): each GPU spreads its links over the
switches, and the switches connect every pair of GPUs with higher
bandwidth and lower latency than a chain of direct links could.

\margin{Lesson~5's unified virtual addressing is what lets a kernel on
one GPU dereference a pointer into a peer GPU's memory, once
\texttt{cudaDeviceEnablePeerAccess} has enabled peer-to-peer access
over PCIe or NVLink (CPG §3.2.9.4 Peer-to-Peer Memory Access).}%
Within a node, then, the GPUs form one NUMA machine: one address space,
with near and far memory.

\subsection{Across boards: RDMA}

GPUs on different boards or in different nodes are attached to their
hosts through PCIe (\cref{fig:l09-nodes}), and PCIe through the CPU
does not provide enough bandwidth to share memory between them.  Their
memories are separate spaces, and the system is a distributed-memory
system as in Lesson~2: data moves by explicit messages, here by RDMA.\margin{Programs rarely issue these transfers
themselves: collective libraries such as NCCL (Lesson~5) choose NVLink,
NVSwitch or GPUDirect RDMA for each pair of GPUs.}

\begin{figure}[tb]
  \centering
  % Two multi-GPU nodes (Lesson 9).  Inside each node, a CPU and a network
% interface card (NIC) sit on a PCIe switch, and four GPUs with their
% memories are joined all-to-all by NVLink: one shared, NUMA address
% space.  Between the nodes, the GPU memories are separate (distributed
% memory) and data moves by RDMA through the NICs and the network.
% Redraws outline figure m9l1_mem_space (with the NIC made explicit).
% Usage: % Two multi-GPU nodes (Lesson 9).  Inside each node, a CPU and a network
% interface card (NIC) sit on a PCIe switch, and four GPUs with their
% memories are joined all-to-all by NVLink: one shared, NUMA address
% space.  Between the nodes, the GPU memories are separate (distributed
% memory) and data moves by RDMA through the NICs and the network.
% Redraws outline figure m9l1_mem_space (with the NIC made explicit).
% Usage: % Two multi-GPU nodes (Lesson 9).  Inside each node, a CPU and a network
% interface card (NIC) sit on a PCIe switch, and four GPUs with their
% memories are joined all-to-all by NVLink: one shared, NUMA address
% space.  Between the nodes, the GPU memories are separate (distributed
% memory) and data moves by RDMA through the NICs and the network.
% Redraws outline figure m9l1_mem_space (with the NIC made explicit).
% Usage: \input{figures/l09-nodes}
% Width: body (120 mm); the drawing is about 118 mm wide.
\begin{tikzpicture}[x=1mm, y=1mm, font=\scriptsize\sffamily,
  gpu/.style={draw, fill=ln-accent!30, minimum width=9mm, minimum height=5mm, inner sep=1pt},
  mem/.style={draw, fill=ln-muted!20, minimum width=4mm, minimum height=5mm, inner sep=0pt},
  box/.style={draw, fill=ln-box, minimum width=9mm, minimum height=5mm, inner sep=1pt},
  sw/.style={draw, fill=ln-muted!60, text=white, minimum height=4.5mm, inner sep=1pt},
  nvl/.style={line width=1.2pt, ln-source!70},
  pcie/.style={line width=0.6pt, ln-muted}]
  \foreach \n/\ox/\lab in {0/0/A, 1/64/B} {
    \begin{scope}[shift={(\ox,0)}]
      \draw[ln-rule, rounded corners=2pt] (0,0) rectangle (54,40);
      \node[anchor=north west] at (1,39.5) {\iflangja{ノード}{node} \lab};
      \node[box] (cpu\n) at (14,33) {CPU};
      \node[box] (nic\n) at (40,33) {NIC};
      \node[sw, minimum width=36mm] (ps\n) at (27,25) {\iflangja{PCIe スイッチ}{PCIe switch}};
      \draw[pcie] (cpu\n) -- (cpu\n|-ps\n.north);
      \draw[pcie] (nic\n) -- (nic\n|-ps\n.north);
      \node[gpu] (a\n) at (17,16) {GPU};
      \node[gpu] (b\n) at (37,16) {GPU};
      \node[gpu] (c\n) at (17,5)  {GPU};
      \node[gpu] (d\n) at (37,5)  {GPU};
      \node[mem] at (9,16) {}; \node[mem] at (45,16) {};
      \node[mem] at (9,5) {};  \node[mem] at (45,5) {};
      \draw[nvl] (a\n) -- (b\n); \draw[nvl] (c\n) -- (d\n);
      \draw[nvl] (a\n) -- (c\n); \draw[nvl] (b\n) -- (d\n);
      \draw[nvl] (a\n) -- (d\n); \draw[nvl] (b\n) -- (c\n);
      \draw[pcie] (a\n.north) -- (a\n|-ps\n.south);
      \draw[pcie] (b\n.north) -- (b\n|-ps\n.south);
    \end{scope}
  }
  % network between the NICs
  \draw[{Stealth[length=2mm]}-{Stealth[length=2mm]}, line width=1.2pt, ln-caution]
     (nic0.north) -- ++(0,4) -| node[pos=0.25, above] {\iflangja{ネットワーク（RDMA）}{network (RDMA)}} (nic1.north);
  % memory labels
  \node[text=ln-muted, align=center] at (27,-4) {\iflangja{共有アドレス空間（NUMA）}{one shared address space (NUMA)}};
  \node[text=ln-muted, align=center] at (91,-4) {\iflangja{共有アドレス空間（NUMA）}{one shared address space (NUMA)}};
  \node[text=ln-caution, align=center] at (59,-9) {\iflangja{ノード間：分散メモリ}{between the nodes: distributed memory}};
  % legend
  \draw[nvl] (0,-12) -- (6,-12); \node[anchor=west] at (6,-12) {NVLink};
  \draw[pcie] (20,-12) -- (26,-12); \node[anchor=west] at (26,-12) {PCIe};
  \node[mem] at (92,-12) {}; \node[anchor=west] at (95,-12) {\iflangja{GPU メモリ}{GPU memory}};
\end{tikzpicture}

% Width: body (120 mm); the drawing is about 118 mm wide.
\begin{tikzpicture}[x=1mm, y=1mm, font=\scriptsize\sffamily,
  gpu/.style={draw, fill=ln-accent!30, minimum width=9mm, minimum height=5mm, inner sep=1pt},
  mem/.style={draw, fill=ln-muted!20, minimum width=4mm, minimum height=5mm, inner sep=0pt},
  box/.style={draw, fill=ln-box, minimum width=9mm, minimum height=5mm, inner sep=1pt},
  sw/.style={draw, fill=ln-muted!60, text=white, minimum height=4.5mm, inner sep=1pt},
  nvl/.style={line width=1.2pt, ln-source!70},
  pcie/.style={line width=0.6pt, ln-muted}]
  \foreach \n/\ox/\lab in {0/0/A, 1/64/B} {
    \begin{scope}[shift={(\ox,0)}]
      \draw[ln-rule, rounded corners=2pt] (0,0) rectangle (54,40);
      \node[anchor=north west] at (1,39.5) {\iflangja{ノード}{node} \lab};
      \node[box] (cpu\n) at (14,33) {CPU};
      \node[box] (nic\n) at (40,33) {NIC};
      \node[sw, minimum width=36mm] (ps\n) at (27,25) {\iflangja{PCIe スイッチ}{PCIe switch}};
      \draw[pcie] (cpu\n) -- (cpu\n|-ps\n.north);
      \draw[pcie] (nic\n) -- (nic\n|-ps\n.north);
      \node[gpu] (a\n) at (17,16) {GPU};
      \node[gpu] (b\n) at (37,16) {GPU};
      \node[gpu] (c\n) at (17,5)  {GPU};
      \node[gpu] (d\n) at (37,5)  {GPU};
      \node[mem] at (9,16) {}; \node[mem] at (45,16) {};
      \node[mem] at (9,5) {};  \node[mem] at (45,5) {};
      \draw[nvl] (a\n) -- (b\n); \draw[nvl] (c\n) -- (d\n);
      \draw[nvl] (a\n) -- (c\n); \draw[nvl] (b\n) -- (d\n);
      \draw[nvl] (a\n) -- (d\n); \draw[nvl] (b\n) -- (c\n);
      \draw[pcie] (a\n.north) -- (a\n|-ps\n.south);
      \draw[pcie] (b\n.north) -- (b\n|-ps\n.south);
    \end{scope}
  }
  % network between the NICs
  \draw[{Stealth[length=2mm]}-{Stealth[length=2mm]}, line width=1.2pt, ln-caution]
     (nic0.north) -- ++(0,4) -| node[pos=0.25, above] {\iflangja{ネットワーク（RDMA）}{network (RDMA)}} (nic1.north);
  % memory labels
  \node[text=ln-muted, align=center] at (27,-4) {\iflangja{共有アドレス空間（NUMA）}{one shared address space (NUMA)}};
  \node[text=ln-muted, align=center] at (91,-4) {\iflangja{共有アドレス空間（NUMA）}{one shared address space (NUMA)}};
  \node[text=ln-caution, align=center] at (59,-9) {\iflangja{ノード間：分散メモリ}{between the nodes: distributed memory}};
  % legend
  \draw[nvl] (0,-12) -- (6,-12); \node[anchor=west] at (6,-12) {NVLink};
  \draw[pcie] (20,-12) -- (26,-12); \node[anchor=west] at (26,-12) {PCIe};
  \node[mem] at (92,-12) {}; \node[anchor=west] at (95,-12) {\iflangja{GPU メモリ}{GPU memory}};
\end{tikzpicture}

% Width: body (120 mm); the drawing is about 118 mm wide.
\begin{tikzpicture}[x=1mm, y=1mm, font=\scriptsize\sffamily,
  gpu/.style={draw, fill=ln-accent!30, minimum width=9mm, minimum height=5mm, inner sep=1pt},
  mem/.style={draw, fill=ln-muted!20, minimum width=4mm, minimum height=5mm, inner sep=0pt},
  box/.style={draw, fill=ln-box, minimum width=9mm, minimum height=5mm, inner sep=1pt},
  sw/.style={draw, fill=ln-muted!60, text=white, minimum height=4.5mm, inner sep=1pt},
  nvl/.style={line width=1.2pt, ln-source!70},
  pcie/.style={line width=0.6pt, ln-muted}]
  \foreach \n/\ox/\lab in {0/0/A, 1/64/B} {
    \begin{scope}[shift={(\ox,0)}]
      \draw[ln-rule, rounded corners=2pt] (0,0) rectangle (54,40);
      \node[anchor=north west] at (1,39.5) {\iflangja{ノード}{node} \lab};
      \node[box] (cpu\n) at (14,33) {CPU};
      \node[box] (nic\n) at (40,33) {NIC};
      \node[sw, minimum width=36mm] (ps\n) at (27,25) {\iflangja{PCIe スイッチ}{PCIe switch}};
      \draw[pcie] (cpu\n) -- (cpu\n|-ps\n.north);
      \draw[pcie] (nic\n) -- (nic\n|-ps\n.north);
      \node[gpu] (a\n) at (17,16) {GPU};
      \node[gpu] (b\n) at (37,16) {GPU};
      \node[gpu] (c\n) at (17,5)  {GPU};
      \node[gpu] (d\n) at (37,5)  {GPU};
      \node[mem] at (9,16) {}; \node[mem] at (45,16) {};
      \node[mem] at (9,5) {};  \node[mem] at (45,5) {};
      \draw[nvl] (a\n) -- (b\n); \draw[nvl] (c\n) -- (d\n);
      \draw[nvl] (a\n) -- (c\n); \draw[nvl] (b\n) -- (d\n);
      \draw[nvl] (a\n) -- (d\n); \draw[nvl] (b\n) -- (c\n);
      \draw[pcie] (a\n.north) -- (a\n|-ps\n.south);
      \draw[pcie] (b\n.north) -- (b\n|-ps\n.south);
    \end{scope}
  }
  % network between the NICs
  \draw[{Stealth[length=2mm]}-{Stealth[length=2mm]}, line width=1.2pt, ln-caution]
     (nic0.north) -- ++(0,4) -| node[pos=0.25, above] {\iflangja{ネットワーク（RDMA）}{network (RDMA)}} (nic1.north);
  % memory labels
  \node[text=ln-muted, align=center] at (27,-4) {\iflangja{共有アドレス空間（NUMA）}{one shared address space (NUMA)}};
  \node[text=ln-muted, align=center] at (91,-4) {\iflangja{共有アドレス空間（NUMA）}{one shared address space (NUMA)}};
  \node[text=ln-caution, align=center] at (59,-9) {\iflangja{ノード間：分散メモリ}{between the nodes: distributed memory}};
  % legend
  \draw[nvl] (0,-12) -- (6,-12); \node[anchor=west] at (6,-12) {NVLink};
  \draw[pcie] (20,-12) -- (26,-12); \node[anchor=west] at (26,-12) {PCIe};
  \node[mem] at (92,-12) {}; \node[anchor=west] at (95,-12) {\iflangja{GPU メモリ}{GPU memory}};
\end{tikzpicture}

  \caption{Two multi-GPU nodes.  Inside a node, NVLink joins the four
    GPUs into one NUMA address space; the CPU and the network interface
    card sit on the PCIe switch (only two GPU connections drawn).  Between
    nodes the memories are separate, and data crosses the network by
    RDMA.}
  \label{fig:l09-nodes}
\end{figure}

Traditional networking sends a buffer through the TCP/IP software stack
(\cref{fig:l09-rdma}a): the sending CPU copies the data from the
application buffer through the operating system to the network, and the
receiving CPU copies it up again into the destination buffer.  Both
CPUs spend cycles on every message, and a busy CPU delays the transfer.

\begin{definition}[RDMA, host bypass, NIC]
  A \term{network interface card (NIC)}[ネットワークインタフェースカード]
  is the device that connects a host to the network.
  The technique called
  \term{remote direct memory access (RDMA)}[リモートダイレクトメモリアクセス（RDMA）]
  extends the DMA of Lesson~7 across the network: the RDMA engine of
  one NIC reads the source buffer from memory and the RDMA engine of the
  other NIC writes it into the destination memory, without either CPU
  copying the data (\cref{fig:l09-rdma}b).  Such a transfer is a
  \term{host bypass}[ホストバイパス] technique: the host processors and
  their software stacks are left out of the data path.
  \index{RDMA|see{remote direct memory access (RDMA)}}%
  \index{NIC|see{network interface card (NIC)}}%
\end{definition}

\begin{figure}[tb]
  \centering
  % Three ways to move a buffer between two hosts (Lesson 9).  (a)
% Traditional networking: the data passes through both CPUs and their
% network stacks.  (b) RDMA: the RDMA engines of the two NICs move the
% data memory to memory; the CPUs are not involved (host bypass).  (c)
% GPUDirect RDMA: the NIC reads and writes GPU memory directly over PCIe,
% so a GPU buffer reaches a remote GPU (or another device) without a copy
% through host memory.
% Redraws outline figures m9l1_rdma, m9l1_rdma1 and m9l1_rdma2.
% Usage: % Three ways to move a buffer between two hosts (Lesson 9).  (a)
% Traditional networking: the data passes through both CPUs and their
% network stacks.  (b) RDMA: the RDMA engines of the two NICs move the
% data memory to memory; the CPUs are not involved (host bypass).  (c)
% GPUDirect RDMA: the NIC reads and writes GPU memory directly over PCIe,
% so a GPU buffer reaches a remote GPU (or another device) without a copy
% through host memory.
% Redraws outline figures m9l1_rdma, m9l1_rdma1 and m9l1_rdma2.
% Usage: % Three ways to move a buffer between two hosts (Lesson 9).  (a)
% Traditional networking: the data passes through both CPUs and their
% network stacks.  (b) RDMA: the RDMA engines of the two NICs move the
% data memory to memory; the CPUs are not involved (host bypass).  (c)
% GPUDirect RDMA: the NIC reads and writes GPU memory directly over PCIe,
% so a GPU buffer reaches a remote GPU (or another device) without a copy
% through host memory.
% Redraws outline figures m9l1_rdma, m9l1_rdma1 and m9l1_rdma2.
% Usage: \input{figures/l09-rdma}
% Width: body (120 mm); the drawing is about 116 mm wide.
\begin{tikzpicture}[x=1mm, y=1mm, font=\scriptsize\sffamily,
  b/.style={draw, fill=ln-box, minimum width=15mm, minimum height=5.5mm, inner sep=1pt, align=center},
  idle/.style={draw=ln-rule, text=ln-muted, fill=white, minimum width=11mm, minimum height=5.5mm, inner sep=1pt, dashed},
  eng/.style={draw, fill=ln-tip!35, minimum width=17mm, minimum height=5.5mm, inner sep=1pt, align=center},
  gpu/.style={draw, fill=ln-accent!30, minimum width=15mm, minimum height=5.5mm, inner sep=1pt, align=center},
  net/.style={draw, fill=ln-caution!20, minimum width=14mm, minimum height=5.5mm, inner sep=1pt},
  a/.style={-{Stealth[length=1.8mm]}, line width=0.7pt},
  host/.style={draw=ln-rule, rounded corners=2pt}]
  % row (a)
  \draw[host] (4,31) rectangle (49,41); \draw[host] (69,31) rectangle (114,41);
  \node[b] (a1) at (13,36) {\iflangja{バッファ 1}{buffer 1}};
  \node[b, fill=ln-muted!20] (a2) at (38,36) {CPU1};
  \node[net] (a3) at (59,36) {\iflangja{ネットワーク}{network}};
  \node[b, fill=ln-muted!20] (a4) at (80,36) {CPU2};
  \node[b] (a5) at (105,36) {\iflangja{バッファ 2}{buffer 2}};
  \draw[a] (a1) -- (a2); \draw[a] (a2) -- (a3); \draw[a] (a3) -- (a4); \draw[a] (a4) -- (a5);
  \node at (0,36) {(a)};
  % row (b)
  \draw[host] (4,15) rectangle (49,28); \draw[host] (69,15) rectangle (114,28);
  \node[b] (b1) at (13,19) {\iflangja{バッファ 1}{buffer 1}};
  \node[idle] at (25,24.5) {CPU1};
  \node[eng] (b2) at (38,19) {\iflangja{RDMA エンジン}{RDMA engine}};
  \node[net] (b3) at (59,19) {\iflangja{ネットワーク}{network}};
  \node[eng] (b4) at (80,19) {\iflangja{RDMA エンジン}{RDMA engine}};
  \node[idle] at (93,24.5) {CPU2};
  \node[b] (b5) at (105,19) {\iflangja{バッファ 2}{buffer 2}};
  \draw[a] (b1) -- (b2); \draw[a] (b2) -- (b3); \draw[a] (b3) -- (b4); \draw[a] (b4) -- (b5);
  \node at (0,21.5) {(b)};
  % row (c)
  \draw[host] (4,0) rectangle (49,13); \draw[host] (69,0) rectangle (114,13);
  \node[gpu] (c1) at (13,4) {GPU1\\\iflangja{メモリ}{memory}};
  \node[idle] at (25,9.5) {CPU1};
  \node[eng] (c2) at (38,4) {NIC};
  \node[net] (c3) at (59,4) {\iflangja{ネットワーク}{network}};
  \node[eng] (c4) at (80,4) {NIC};
  \node[idle] at (93,9.5) {CPU2};
  \node[gpu] (c5) at (105,4) {GPU2\\\iflangja{メモリ}{memory}};
  \draw[a] (c1) -- (c2); \draw[a] (c2) -- (c3); \draw[a] (c3) -- (c4); \draw[a] (c4) -- (c5);
  \node at (0,6.5) {(c)};
\end{tikzpicture}

% Width: body (120 mm); the drawing is about 116 mm wide.
\begin{tikzpicture}[x=1mm, y=1mm, font=\scriptsize\sffamily,
  b/.style={draw, fill=ln-box, minimum width=15mm, minimum height=5.5mm, inner sep=1pt, align=center},
  idle/.style={draw=ln-rule, text=ln-muted, fill=white, minimum width=11mm, minimum height=5.5mm, inner sep=1pt, dashed},
  eng/.style={draw, fill=ln-tip!35, minimum width=17mm, minimum height=5.5mm, inner sep=1pt, align=center},
  gpu/.style={draw, fill=ln-accent!30, minimum width=15mm, minimum height=5.5mm, inner sep=1pt, align=center},
  net/.style={draw, fill=ln-caution!20, minimum width=14mm, minimum height=5.5mm, inner sep=1pt},
  a/.style={-{Stealth[length=1.8mm]}, line width=0.7pt},
  host/.style={draw=ln-rule, rounded corners=2pt}]
  % row (a)
  \draw[host] (4,31) rectangle (49,41); \draw[host] (69,31) rectangle (114,41);
  \node[b] (a1) at (13,36) {\iflangja{バッファ 1}{buffer 1}};
  \node[b, fill=ln-muted!20] (a2) at (38,36) {CPU1};
  \node[net] (a3) at (59,36) {\iflangja{ネットワーク}{network}};
  \node[b, fill=ln-muted!20] (a4) at (80,36) {CPU2};
  \node[b] (a5) at (105,36) {\iflangja{バッファ 2}{buffer 2}};
  \draw[a] (a1) -- (a2); \draw[a] (a2) -- (a3); \draw[a] (a3) -- (a4); \draw[a] (a4) -- (a5);
  \node at (0,36) {(a)};
  % row (b)
  \draw[host] (4,15) rectangle (49,28); \draw[host] (69,15) rectangle (114,28);
  \node[b] (b1) at (13,19) {\iflangja{バッファ 1}{buffer 1}};
  \node[idle] at (25,24.5) {CPU1};
  \node[eng] (b2) at (38,19) {\iflangja{RDMA エンジン}{RDMA engine}};
  \node[net] (b3) at (59,19) {\iflangja{ネットワーク}{network}};
  \node[eng] (b4) at (80,19) {\iflangja{RDMA エンジン}{RDMA engine}};
  \node[idle] at (93,24.5) {CPU2};
  \node[b] (b5) at (105,19) {\iflangja{バッファ 2}{buffer 2}};
  \draw[a] (b1) -- (b2); \draw[a] (b2) -- (b3); \draw[a] (b3) -- (b4); \draw[a] (b4) -- (b5);
  \node at (0,21.5) {(b)};
  % row (c)
  \draw[host] (4,0) rectangle (49,13); \draw[host] (69,0) rectangle (114,13);
  \node[gpu] (c1) at (13,4) {GPU1\\\iflangja{メモリ}{memory}};
  \node[idle] at (25,9.5) {CPU1};
  \node[eng] (c2) at (38,4) {NIC};
  \node[net] (c3) at (59,4) {\iflangja{ネットワーク}{network}};
  \node[eng] (c4) at (80,4) {NIC};
  \node[idle] at (93,9.5) {CPU2};
  \node[gpu] (c5) at (105,4) {GPU2\\\iflangja{メモリ}{memory}};
  \draw[a] (c1) -- (c2); \draw[a] (c2) -- (c3); \draw[a] (c3) -- (c4); \draw[a] (c4) -- (c5);
  \node at (0,6.5) {(c)};
\end{tikzpicture}

% Width: body (120 mm); the drawing is about 116 mm wide.
\begin{tikzpicture}[x=1mm, y=1mm, font=\scriptsize\sffamily,
  b/.style={draw, fill=ln-box, minimum width=15mm, minimum height=5.5mm, inner sep=1pt, align=center},
  idle/.style={draw=ln-rule, text=ln-muted, fill=white, minimum width=11mm, minimum height=5.5mm, inner sep=1pt, dashed},
  eng/.style={draw, fill=ln-tip!35, minimum width=17mm, minimum height=5.5mm, inner sep=1pt, align=center},
  gpu/.style={draw, fill=ln-accent!30, minimum width=15mm, minimum height=5.5mm, inner sep=1pt, align=center},
  net/.style={draw, fill=ln-caution!20, minimum width=14mm, minimum height=5.5mm, inner sep=1pt},
  a/.style={-{Stealth[length=1.8mm]}, line width=0.7pt},
  host/.style={draw=ln-rule, rounded corners=2pt}]
  % row (a)
  \draw[host] (4,31) rectangle (49,41); \draw[host] (69,31) rectangle (114,41);
  \node[b] (a1) at (13,36) {\iflangja{バッファ 1}{buffer 1}};
  \node[b, fill=ln-muted!20] (a2) at (38,36) {CPU1};
  \node[net] (a3) at (59,36) {\iflangja{ネットワーク}{network}};
  \node[b, fill=ln-muted!20] (a4) at (80,36) {CPU2};
  \node[b] (a5) at (105,36) {\iflangja{バッファ 2}{buffer 2}};
  \draw[a] (a1) -- (a2); \draw[a] (a2) -- (a3); \draw[a] (a3) -- (a4); \draw[a] (a4) -- (a5);
  \node at (0,36) {(a)};
  % row (b)
  \draw[host] (4,15) rectangle (49,28); \draw[host] (69,15) rectangle (114,28);
  \node[b] (b1) at (13,19) {\iflangja{バッファ 1}{buffer 1}};
  \node[idle] at (25,24.5) {CPU1};
  \node[eng] (b2) at (38,19) {\iflangja{RDMA エンジン}{RDMA engine}};
  \node[net] (b3) at (59,19) {\iflangja{ネットワーク}{network}};
  \node[eng] (b4) at (80,19) {\iflangja{RDMA エンジン}{RDMA engine}};
  \node[idle] at (93,24.5) {CPU2};
  \node[b] (b5) at (105,19) {\iflangja{バッファ 2}{buffer 2}};
  \draw[a] (b1) -- (b2); \draw[a] (b2) -- (b3); \draw[a] (b3) -- (b4); \draw[a] (b4) -- (b5);
  \node at (0,21.5) {(b)};
  % row (c)
  \draw[host] (4,0) rectangle (49,13); \draw[host] (69,0) rectangle (114,13);
  \node[gpu] (c1) at (13,4) {GPU1\\\iflangja{メモリ}{memory}};
  \node[idle] at (25,9.5) {CPU1};
  \node[eng] (c2) at (38,4) {NIC};
  \node[net] (c3) at (59,4) {\iflangja{ネットワーク}{network}};
  \node[eng] (c4) at (80,4) {NIC};
  \node[idle] at (93,9.5) {CPU2};
  \node[gpu] (c5) at (105,4) {GPU2\\\iflangja{メモリ}{memory}};
  \draw[a] (c1) -- (c2); \draw[a] (c2) -- (c3); \draw[a] (c3) -- (c4); \draw[a] (c4) -- (c5);
  \node at (0,6.5) {(c)};
\end{tikzpicture}

  \caption{(a) Traditional networking: the data passes through both CPUs.
    (b) RDMA: the RDMA engines move it memory to memory; the CPUs
    (dashed) stay out of the path.  (c) GPUDirect RDMA: the NIC reads and
    writes GPU memory directly, so a GPU buffer reaches a remote GPU
    without a copy through host memory.}
  \label{fig:l09-rdma}
\end{figure}

\begin{definition}[GPUDirect RDMA]
  NVIDIA's technology
  \term{GPUDirect RDMA}[GPUDirect RDMA] provides a direct data path
  between GPU memory and a third-party PCIe device, such as a NIC or a
  storage adapter, \enquote{using standard features of PCI Express}
  \cite{nvidia2026gpudirectrdma}.  A GPU buffer can then
  be sent through the network to a remote GPU, or to another device,
  without being staged in host memory (\cref{fig:l09-rdma}c).
\end{definition}

\caution{GPUDirect RDMA needs the GPU and the third-party device to
\enquote{share the same upstream PCI Express root complex}; it was
introduced with Kepler-class GPUs and CUDA~5.0
\cite{nvidia2026gpudirectrdma}.}%
The two kinds of link work together: inside a node, NVLink and NVSwitch
carry the GPU-to-GPU traffic; between nodes, GPUDirect RDMA moves it
over the network.  PCIe alone cannot meet the demand of so many GPUs
exchanging data.

\begin{keypoint}
  Scaling past one die leads to multi-chip modules, multi-GPU nodes and
  clusters.  Within a package or node the memory is one NUMA space, so
  block scheduling and page placement decide how many accesses are far;
  NVLink and NVSwitch make far accesses cheaper.  Between nodes the
  memory is distributed and RDMA, with GPUDirect RDMA for GPU memory,
  moves the data without the CPUs.
\end{keypoint}

\section{Concurrency on GPUs}
\label{sec:l09-concurrency}

A large GPU is not always used well.\source{M9, concurrency on GPUs;
CPG §3.2.8 Asynchronous Concurrent Execution, §3.2.9 Multi-Device
System; \cite{nvidia2020ampere}; \cite{nvidia2026mps}.}  This section
asks how several jobs can share one GPU, compares three mechanisms that
let them, and shows how a program drives several streams and several
GPUs.

\subsection{Utilization and multi-tenancy}

\begin{definition}[Utilization, datacenter, multi-tenancy]
  The \term{utilization}[利用率] of a GPU is the fraction of its
  resources (SMs, memory bandwidth, time) that does useful work.  A
  \term{datacenter}[データセンタ] is a facility of many servers
  shared by many users, and \term{multi-tenancy}[マルチテナンシ] is the
  sharing of one machine, here one GPU, by the jobs of several
  independent users (tenants).
\end{definition}

Not every job can use all of a large GPU: an inference service with
small batches, a development session or a small training job may keep a
fraction of the SMs busy.  On a dedicated GPU the rest is wasted, which
matters most in datacenters, where GPUs are the expensive resource.
Traditional datacenters serve many tenants with jobs of this size.  Some
AI workloads behave differently: training a large language model takes
whole GPUs, many of them, for a single tenant, and such a job advances
only as fast as its slowest task, since every step waits for all of its
GPUs.\margin{Ye et al.\ survey the scheduling of deep-learning jobs in
GPU datacenters and note that \enquote{traditional approaches designed
for big data or high performance computing workloads can not support DL
workloads to fully utilize the GPU resources} \cite{ye2024dlsched}; Yu et
al.\ survey multi-tenant inference \cite{yu2022multitenant}.}  Both
patterns call for ways to place more than one stream of work on a GPU.

\subsection{Three mechanisms}

\Cref{fig:l09-concurrency} compares the three mechanisms by what they
partition.

\begin{figure}[tb]
  \centering
  % GPU concurrency mechanisms (Lesson 9).  (a) One kernel: an SM with
% its L1, the shared L2 and DRAM.  (b) Streams of one client: work of
% streams A, B and C mixes on every SM, in the L2 and in DRAM; nothing is
% partitioned.  (c) MPS: the clients' blocks are spread over the SMs,
% drawn with a per-client SM share as in the lecture (MPS limits a
% client's share but reserves no SMs); the L2 and DRAM are shared.
% (d) MIG: SMs, L2 slices and DRAM (bandwidth and capacity) are split
% physically between the clients (dashed line).
% Redraws outline figure m9l2_gpucon.
% Usage: % GPU concurrency mechanisms (Lesson 9).  (a) One kernel: an SM with
% its L1, the shared L2 and DRAM.  (b) Streams of one client: work of
% streams A, B and C mixes on every SM, in the L2 and in DRAM; nothing is
% partitioned.  (c) MPS: the clients' blocks are spread over the SMs,
% drawn with a per-client SM share as in the lecture (MPS limits a
% client's share but reserves no SMs); the L2 and DRAM are shared.
% (d) MIG: SMs, L2 slices and DRAM (bandwidth and capacity) are split
% physically between the clients (dashed line).
% Redraws outline figure m9l2_gpucon.
% Usage: % GPU concurrency mechanisms (Lesson 9).  (a) One kernel: an SM with
% its L1, the shared L2 and DRAM.  (b) Streams of one client: work of
% streams A, B and C mixes on every SM, in the L2 and in DRAM; nothing is
% partitioned.  (c) MPS: the clients' blocks are spread over the SMs,
% drawn with a per-client SM share as in the lecture (MPS limits a
% client's share but reserves no SMs); the L2 and DRAM are shared.
% (d) MIG: SMs, L2 slices and DRAM (bandwidth and capacity) are split
% physically between the clients (dashed line).
% Redraws outline figure m9l2_gpucon.
% Usage: \input{figures/l09-concurrency}
% Width: body (120 mm); the drawing is about 116 mm wide.
\begin{tikzpicture}[x=1mm, y=1mm, font=\scriptsize\sffamily,
  tile/.style={draw, minimum width=5.4mm, minimum height=4.2mm, inner sep=0pt},
  inp/.style={draw, minimum height=4.5mm, minimum width=6mm, inner sep=1pt, align=center},
  sch/.style={draw, fill=white, minimum width=24mm, minimum height=4mm, inner sep=0pt},
  bar/.style={draw, fill=white, minimum width=24mm, minimum height=4mm, inner sep=0pt},
  a/.style={-{Stealth[length=1.5mm]}, line width=0.5pt}]
  \colorlet{cA}{ln-accent!35}\colorlet{cB}{ln-tip!40}\colorlet{cC}{ln-caution!30}
  % \stripes{x}{y}{w}{h}: rectangle filled with stripes of the three colours
  \newcommand{\stripes}[4]{%
    \pgfmathtruncatemacro{\n}{floor((#3)/1.2)}%
    \foreach \s in {0,...,\n} {%
      \pgfmathsetmacro{\xx}{#1+\s*1.2}%
      \pgfmathsetmacro{\ww}{min(1.2,#1+#3-\xx)}%
      \pgfmathtruncatemacro{\c}{mod(\s,3)}%
      \ifdim\ww pt>0pt
        \ifnum\c=0 \fill[cA] (\xx,#2) rectangle ++(\ww,#4);\fi
        \ifnum\c=1 \fill[cB] (\xx,#2) rectangle ++(\ww,#4);\fi
        \ifnum\c=2 \fill[cC] (\xx,#2) rectangle ++(\ww,#4);\fi
      \fi
    }%
    \draw (#1,#2) rectangle ++(#3,#4);}
  % ------------------------------------------------ (a) one kernel
  \node at (13,49) {\iflangja{(a) 単一カーネル}{(a) one kernel}};
  \node[inp, fill=ln-box, minimum width=14mm] (k) at (13,42) {\iflangja{カーネル}{kernel}};
  \node[draw, fill=ln-box, minimum width=18mm, minimum height=11mm] (sm) at (13,24.5) {};
  \node[anchor=north] at (sm.north) {SM};
  \node[draw, fill=white, minimum width=12mm, minimum height=3.5mm, inner sep=0pt] at (13,22.5) {L1};
  \draw[a] (k) -- (sm);
  \node[bar] at (13,14) {L2};
  \node[bar] at (13,7) {DRAM};
  \draw[a] (sm) -- (13,16);
  % ------------------------------------------------ (b) streams
  \begin{scope}[shift={(30,0)}]
    \node at (13,49) {\iflangja{(b) ストリーム}{(b) streams}};
    \node[inp, fill=cA] at (5,42) {A};
    \node[inp, fill=cB] at (13,42) {B};
    \node[inp, fill=cC] at (21,42) {C};
    \node[sch] at (13,35) {\iflangja{スケジューラ}{scheduler}};
    \foreach \x in {5,13,21} \draw[a] (\x,39.7) -- (\x,37);
    \foreach \c in {0,...,3} \foreach \r in {0,1} {
      \stripes{1.35+\c*6.0}{18.5+\r*5.5}{5.4}{4.2}
    }
    \draw[a] (13,33) -- (13,28.7);
    \stripes{1}{12}{24}{4}\node[fill=white, inner sep=0.5pt] at (13,14) {L2};
    \stripes{1}{5}{24}{4}\node[fill=white, inner sep=0.5pt] at (13,7) {DRAM};
  \end{scope}
  % ------------------------------------------------ (c) MPS
  \begin{scope}[shift={(60,0)}]
    \node at (13,49) {(c) MPS};
    \node[inp, fill=cA] at (7,42) {1};
    \node[inp, fill=cB] at (19,42) {2};
    \node[sch] at (13,35) {\iflangja{スケジューラ}{scheduler}};
    \foreach \x in {7,19} \draw[a] (\x,39.7) -- (\x,37);
    \foreach \c/\r/\col in {0/1/cA,1/1/cA,2/1/cB,3/1/cA,0/0/cB,1/0/cB,2/0/cA,3/0/cB} {
      \node[tile, fill=\col] at (4.05+\c*6.0, 20.6+\r*5.5) {};
    }
    \draw[a] (13,33) -- (13,28.7);
    \node[bar] at (13,14) {\iflangja{L2（共有）}{L2 (shared)}};
    \node[bar] at (13,7) {\iflangja{DRAM（共有）}{DRAM (shared)}};
  \end{scope}
  % ------------------------------------------------ (d) MIG
  \begin{scope}[shift={(90,0)}]
    \node at (13,49) {(d) MIG};
    \node[inp, fill=cA] at (7,42) {1};
    \node[inp, fill=cB] at (19,42) {2};
    \node[sch] at (13,35) {\iflangja{スケジューラ}{scheduler}};
    \foreach \x in {7,19} \draw[a] (\x,39.7) -- (\x,37);
    \foreach \c in {0,1} \foreach \r in {0,1} {
      \node[tile, fill=cA] at (4.05+\c*6.0, 20.6+\r*5.5) {};
      \node[tile, fill=cB] at (16.05+\c*6.0, 20.6+\r*5.5) {};
    }
    \foreach \x in {7,19} \draw[a] (\x,33) -- (\x,28.7);
    \fill[cA] (1,12) rectangle (13,16); \fill[cB] (13,12) rectangle (25,16);
    \draw (1,12) rectangle (25,16); \node[fill=white, inner sep=0.5pt] at (13,14) {L2};
    \fill[cA] (1,5) rectangle (13,9); \fill[cB] (13,5) rectangle (25,9);
    \draw (1,5) rectangle (25,9); \node[fill=white, inner sep=0.5pt] at (13,7) {DRAM};
    \draw[ln-caution, line width=1pt, dashed] (13,45.5) -- (13,37.3) (13,32.7) -- (13,16.3) (13,11.7) -- (13,9.3) (13,4.7) -- (13,3.5);
  \end{scope}
  % legend
  \foreach \col/\lab/\x in {cA/{\iflangja{クライアント 1／ストリーム A}{client 1 / stream A}}/1,
                            cB/{\iflangja{クライアント 2／ストリーム B}{client 2 / stream B}}/42,
                            cC/{\iflangja{ストリーム C}{stream C}}/83} {
    \fill[\col] (\x,-2) rectangle ++(3.5,3); \draw (\x,-2) rectangle ++(3.5,3);
    \node[anchor=west] at (\x+3.5,-0.5) {\lab};
  }
\end{tikzpicture}

% Width: body (120 mm); the drawing is about 116 mm wide.
\begin{tikzpicture}[x=1mm, y=1mm, font=\scriptsize\sffamily,
  tile/.style={draw, minimum width=5.4mm, minimum height=4.2mm, inner sep=0pt},
  inp/.style={draw, minimum height=4.5mm, minimum width=6mm, inner sep=1pt, align=center},
  sch/.style={draw, fill=white, minimum width=24mm, minimum height=4mm, inner sep=0pt},
  bar/.style={draw, fill=white, minimum width=24mm, minimum height=4mm, inner sep=0pt},
  a/.style={-{Stealth[length=1.5mm]}, line width=0.5pt}]
  \colorlet{cA}{ln-accent!35}\colorlet{cB}{ln-tip!40}\colorlet{cC}{ln-caution!30}
  % \stripes{x}{y}{w}{h}: rectangle filled with stripes of the three colours
  \newcommand{\stripes}[4]{%
    \pgfmathtruncatemacro{\n}{floor((#3)/1.2)}%
    \foreach \s in {0,...,\n} {%
      \pgfmathsetmacro{\xx}{#1+\s*1.2}%
      \pgfmathsetmacro{\ww}{min(1.2,#1+#3-\xx)}%
      \pgfmathtruncatemacro{\c}{mod(\s,3)}%
      \ifdim\ww pt>0pt
        \ifnum\c=0 \fill[cA] (\xx,#2) rectangle ++(\ww,#4);\fi
        \ifnum\c=1 \fill[cB] (\xx,#2) rectangle ++(\ww,#4);\fi
        \ifnum\c=2 \fill[cC] (\xx,#2) rectangle ++(\ww,#4);\fi
      \fi
    }%
    \draw (#1,#2) rectangle ++(#3,#4);}
  % ------------------------------------------------ (a) one kernel
  \node at (13,49) {\iflangja{(a) 単一カーネル}{(a) one kernel}};
  \node[inp, fill=ln-box, minimum width=14mm] (k) at (13,42) {\iflangja{カーネル}{kernel}};
  \node[draw, fill=ln-box, minimum width=18mm, minimum height=11mm] (sm) at (13,24.5) {};
  \node[anchor=north] at (sm.north) {SM};
  \node[draw, fill=white, minimum width=12mm, minimum height=3.5mm, inner sep=0pt] at (13,22.5) {L1};
  \draw[a] (k) -- (sm);
  \node[bar] at (13,14) {L2};
  \node[bar] at (13,7) {DRAM};
  \draw[a] (sm) -- (13,16);
  % ------------------------------------------------ (b) streams
  \begin{scope}[shift={(30,0)}]
    \node at (13,49) {\iflangja{(b) ストリーム}{(b) streams}};
    \node[inp, fill=cA] at (5,42) {A};
    \node[inp, fill=cB] at (13,42) {B};
    \node[inp, fill=cC] at (21,42) {C};
    \node[sch] at (13,35) {\iflangja{スケジューラ}{scheduler}};
    \foreach \x in {5,13,21} \draw[a] (\x,39.7) -- (\x,37);
    \foreach \c in {0,...,3} \foreach \r in {0,1} {
      \stripes{1.35+\c*6.0}{18.5+\r*5.5}{5.4}{4.2}
    }
    \draw[a] (13,33) -- (13,28.7);
    \stripes{1}{12}{24}{4}\node[fill=white, inner sep=0.5pt] at (13,14) {L2};
    \stripes{1}{5}{24}{4}\node[fill=white, inner sep=0.5pt] at (13,7) {DRAM};
  \end{scope}
  % ------------------------------------------------ (c) MPS
  \begin{scope}[shift={(60,0)}]
    \node at (13,49) {(c) MPS};
    \node[inp, fill=cA] at (7,42) {1};
    \node[inp, fill=cB] at (19,42) {2};
    \node[sch] at (13,35) {\iflangja{スケジューラ}{scheduler}};
    \foreach \x in {7,19} \draw[a] (\x,39.7) -- (\x,37);
    \foreach \c/\r/\col in {0/1/cA,1/1/cA,2/1/cB,3/1/cA,0/0/cB,1/0/cB,2/0/cA,3/0/cB} {
      \node[tile, fill=\col] at (4.05+\c*6.0, 20.6+\r*5.5) {};
    }
    \draw[a] (13,33) -- (13,28.7);
    \node[bar] at (13,14) {\iflangja{L2（共有）}{L2 (shared)}};
    \node[bar] at (13,7) {\iflangja{DRAM（共有）}{DRAM (shared)}};
  \end{scope}
  % ------------------------------------------------ (d) MIG
  \begin{scope}[shift={(90,0)}]
    \node at (13,49) {(d) MIG};
    \node[inp, fill=cA] at (7,42) {1};
    \node[inp, fill=cB] at (19,42) {2};
    \node[sch] at (13,35) {\iflangja{スケジューラ}{scheduler}};
    \foreach \x in {7,19} \draw[a] (\x,39.7) -- (\x,37);
    \foreach \c in {0,1} \foreach \r in {0,1} {
      \node[tile, fill=cA] at (4.05+\c*6.0, 20.6+\r*5.5) {};
      \node[tile, fill=cB] at (16.05+\c*6.0, 20.6+\r*5.5) {};
    }
    \foreach \x in {7,19} \draw[a] (\x,33) -- (\x,28.7);
    \fill[cA] (1,12) rectangle (13,16); \fill[cB] (13,12) rectangle (25,16);
    \draw (1,12) rectangle (25,16); \node[fill=white, inner sep=0.5pt] at (13,14) {L2};
    \fill[cA] (1,5) rectangle (13,9); \fill[cB] (13,5) rectangle (25,9);
    \draw (1,5) rectangle (25,9); \node[fill=white, inner sep=0.5pt] at (13,7) {DRAM};
    \draw[ln-caution, line width=1pt, dashed] (13,45.5) -- (13,37.3) (13,32.7) -- (13,16.3) (13,11.7) -- (13,9.3) (13,4.7) -- (13,3.5);
  \end{scope}
  % legend
  \foreach \col/\lab/\x in {cA/{\iflangja{クライアント 1／ストリーム A}{client 1 / stream A}}/1,
                            cB/{\iflangja{クライアント 2／ストリーム B}{client 2 / stream B}}/42,
                            cC/{\iflangja{ストリーム C}{stream C}}/83} {
    \fill[\col] (\x,-2) rectangle ++(3.5,3); \draw (\x,-2) rectangle ++(3.5,3);
    \node[anchor=west] at (\x+3.5,-0.5) {\lab};
  }
\end{tikzpicture}

% Width: body (120 mm); the drawing is about 116 mm wide.
\begin{tikzpicture}[x=1mm, y=1mm, font=\scriptsize\sffamily,
  tile/.style={draw, minimum width=5.4mm, minimum height=4.2mm, inner sep=0pt},
  inp/.style={draw, minimum height=4.5mm, minimum width=6mm, inner sep=1pt, align=center},
  sch/.style={draw, fill=white, minimum width=24mm, minimum height=4mm, inner sep=0pt},
  bar/.style={draw, fill=white, minimum width=24mm, minimum height=4mm, inner sep=0pt},
  a/.style={-{Stealth[length=1.5mm]}, line width=0.5pt}]
  \colorlet{cA}{ln-accent!35}\colorlet{cB}{ln-tip!40}\colorlet{cC}{ln-caution!30}
  % \stripes{x}{y}{w}{h}: rectangle filled with stripes of the three colours
  \newcommand{\stripes}[4]{%
    \pgfmathtruncatemacro{\n}{floor((#3)/1.2)}%
    \foreach \s in {0,...,\n} {%
      \pgfmathsetmacro{\xx}{#1+\s*1.2}%
      \pgfmathsetmacro{\ww}{min(1.2,#1+#3-\xx)}%
      \pgfmathtruncatemacro{\c}{mod(\s,3)}%
      \ifdim\ww pt>0pt
        \ifnum\c=0 \fill[cA] (\xx,#2) rectangle ++(\ww,#4);\fi
        \ifnum\c=1 \fill[cB] (\xx,#2) rectangle ++(\ww,#4);\fi
        \ifnum\c=2 \fill[cC] (\xx,#2) rectangle ++(\ww,#4);\fi
      \fi
    }%
    \draw (#1,#2) rectangle ++(#3,#4);}
  % ------------------------------------------------ (a) one kernel
  \node at (13,49) {\iflangja{(a) 単一カーネル}{(a) one kernel}};
  \node[inp, fill=ln-box, minimum width=14mm] (k) at (13,42) {\iflangja{カーネル}{kernel}};
  \node[draw, fill=ln-box, minimum width=18mm, minimum height=11mm] (sm) at (13,24.5) {};
  \node[anchor=north] at (sm.north) {SM};
  \node[draw, fill=white, minimum width=12mm, minimum height=3.5mm, inner sep=0pt] at (13,22.5) {L1};
  \draw[a] (k) -- (sm);
  \node[bar] at (13,14) {L2};
  \node[bar] at (13,7) {DRAM};
  \draw[a] (sm) -- (13,16);
  % ------------------------------------------------ (b) streams
  \begin{scope}[shift={(30,0)}]
    \node at (13,49) {\iflangja{(b) ストリーム}{(b) streams}};
    \node[inp, fill=cA] at (5,42) {A};
    \node[inp, fill=cB] at (13,42) {B};
    \node[inp, fill=cC] at (21,42) {C};
    \node[sch] at (13,35) {\iflangja{スケジューラ}{scheduler}};
    \foreach \x in {5,13,21} \draw[a] (\x,39.7) -- (\x,37);
    \foreach \c in {0,...,3} \foreach \r in {0,1} {
      \stripes{1.35+\c*6.0}{18.5+\r*5.5}{5.4}{4.2}
    }
    \draw[a] (13,33) -- (13,28.7);
    \stripes{1}{12}{24}{4}\node[fill=white, inner sep=0.5pt] at (13,14) {L2};
    \stripes{1}{5}{24}{4}\node[fill=white, inner sep=0.5pt] at (13,7) {DRAM};
  \end{scope}
  % ------------------------------------------------ (c) MPS
  \begin{scope}[shift={(60,0)}]
    \node at (13,49) {(c) MPS};
    \node[inp, fill=cA] at (7,42) {1};
    \node[inp, fill=cB] at (19,42) {2};
    \node[sch] at (13,35) {\iflangja{スケジューラ}{scheduler}};
    \foreach \x in {7,19} \draw[a] (\x,39.7) -- (\x,37);
    \foreach \c/\r/\col in {0/1/cA,1/1/cA,2/1/cB,3/1/cA,0/0/cB,1/0/cB,2/0/cA,3/0/cB} {
      \node[tile, fill=\col] at (4.05+\c*6.0, 20.6+\r*5.5) {};
    }
    \draw[a] (13,33) -- (13,28.7);
    \node[bar] at (13,14) {\iflangja{L2（共有）}{L2 (shared)}};
    \node[bar] at (13,7) {\iflangja{DRAM（共有）}{DRAM (shared)}};
  \end{scope}
  % ------------------------------------------------ (d) MIG
  \begin{scope}[shift={(90,0)}]
    \node at (13,49) {(d) MIG};
    \node[inp, fill=cA] at (7,42) {1};
    \node[inp, fill=cB] at (19,42) {2};
    \node[sch] at (13,35) {\iflangja{スケジューラ}{scheduler}};
    \foreach \x in {7,19} \draw[a] (\x,39.7) -- (\x,37);
    \foreach \c in {0,1} \foreach \r in {0,1} {
      \node[tile, fill=cA] at (4.05+\c*6.0, 20.6+\r*5.5) {};
      \node[tile, fill=cB] at (16.05+\c*6.0, 20.6+\r*5.5) {};
    }
    \foreach \x in {7,19} \draw[a] (\x,33) -- (\x,28.7);
    \fill[cA] (1,12) rectangle (13,16); \fill[cB] (13,12) rectangle (25,16);
    \draw (1,12) rectangle (25,16); \node[fill=white, inner sep=0.5pt] at (13,14) {L2};
    \fill[cA] (1,5) rectangle (13,9); \fill[cB] (13,5) rectangle (25,9);
    \draw (1,5) rectangle (25,9); \node[fill=white, inner sep=0.5pt] at (13,7) {DRAM};
    \draw[ln-caution, line width=1pt, dashed] (13,45.5) -- (13,37.3) (13,32.7) -- (13,16.3) (13,11.7) -- (13,9.3) (13,4.7) -- (13,3.5);
  \end{scope}
  % legend
  \foreach \col/\lab/\x in {cA/{\iflangja{クライアント 1／ストリーム A}{client 1 / stream A}}/1,
                            cB/{\iflangja{クライアント 2／ストリーム B}{client 2 / stream B}}/42,
                            cC/{\iflangja{ストリーム C}{stream C}}/83} {
    \fill[\col] (\x,-2) rectangle ++(3.5,3); \draw (\x,-2) rectangle ++(3.5,3);
    \node[anchor=west] at (\x+3.5,-0.5) {\lab};
  }
\end{tikzpicture}

  \caption{GPU concurrency mechanisms.  (a) One kernel uses the SMs, the
    L2 and DRAM.  (b) Streams of one client: their work mixes on every
    SM and in the L2 and DRAM; nothing is partitioned.  (c) MPS: the
    clients' blocks are spread over the SMs (drawn here with a per-client
    SM share); the L2 and DRAM are shared.  (d)
    MIG: SMs, L2 and DRAM are split physically between the clients.}
  \label{fig:l09-concurrency}
\end{figure}

\begin{description}
  \item[Streams.]  One client (one process) issues work to several CUDA
    streams (Lesson~5).  Only independent streams run concurrently, and
    their kernels share all resources: SMs, L2 and DRAM.
  \item[MPS.]  Several clients, separate processes, run at the same time
    on different SMs, while the L2 cache and DRAM bandwidth remain
    shared.
  \item[MIG.]  The GPU is divided physically: each client receives its
    own SMs, L2 slices and DRAM bandwidth.
\end{description}

\subsection{Multi-Instance GPU}

\begin{definition}[Multi-Instance GPU, GPU instance, QoS]
  The technology called
  \term{Multi-Instance GPU (MIG)}[マルチインスタンスGPU（MIG）] divides one
  GPU into several isolated partitions, each a
  \term{GPU instance}[GPUインスタンス] that behaves like a smaller GPU
  with its own SMs, L2 cache slices, memory controllers and DRAM
  \cite{nvidia2020ampere}.  MIG partitions memory, the L2 cache ports,
  DRAM bandwidth and the on-chip crossbar ports, so that each instance
  has a guaranteed
  \term{quality of service (QoS)}[サービス品質（QoS）]: predictable
  throughput and latency whatever the other instances do.
  \index{MIG|see{Multi-Instance GPU (MIG)}}%
  \index{QoS|see{quality of service (QoS)}}%
\end{definition}

\begin{marginfigure}
  \centering
  % MIG instances on a seven-slice GPU (Lesson 9).  Seven GPU slices, each
% with its compute (GPU) and memory part; four applications own GPU
% instances of 3, 2, 1 and 1 slices.  Own allocation; the outline figure
% shows 1, 1, 2 and 3.  The memory side is simplified to one block per
% slice: the A100 has eight memory slices and a 3-slice instance gets
% four (3g.20gb), see the caption.
% Redraws outline figure m9l2_gpumulti (allocation changed).
% Usage: % MIG instances on a seven-slice GPU (Lesson 9).  Seven GPU slices, each
% with its compute (GPU) and memory part; four applications own GPU
% instances of 3, 2, 1 and 1 slices.  Own allocation; the outline figure
% shows 1, 1, 2 and 3.  The memory side is simplified to one block per
% slice: the A100 has eight memory slices and a 3-slice instance gets
% four (3g.20gb), see the caption.
% Redraws outline figure m9l2_gpumulti (allocation changed).
% Usage: % MIG instances on a seven-slice GPU (Lesson 9).  Seven GPU slices, each
% with its compute (GPU) and memory part; four applications own GPU
% instances of 3, 2, 1 and 1 slices.  Own allocation; the outline figure
% shows 1, 1, 2 and 3.  The memory side is simplified to one block per
% slice: the A100 has eight memory slices and a 3-slice instance gets
% four (3g.20gb), see the caption.
% Redraws outline figure m9l2_gpumulti (allocation changed).
% Usage: \input{figures/l09-mig}
% Width: margin (50 mm); the drawing is about 48 mm wide.
\begin{tikzpicture}[x=1mm, y=1mm, font=\scriptsize\sffamily,
  g/.style={draw, fill=ln-accent!30, minimum width=6mm, minimum height=5mm, inner sep=0pt},
  m/.style={draw, fill=ln-tip!35, minimum width=6mm, minimum height=3.5mm, inner sep=0pt},
  app/.style={draw, fill=ln-box, minimum height=4.5mm, inner sep=0pt}]
  \foreach \i in {0,...,6} {
    \node[g] at (3+\i*7,10) {\the\numexpr\i+1\relax};
    \node[m] at (3+\i*7,5) {};
  }
  \node[app, minimum width=20mm] at (10,-1) {App 1};
  \node[app, minimum width=13mm] at (27.5,-1) {App 2};
  \node[app, minimum width=6mm] at (38,-1) {3};
  \node[app, minimum width=6mm] at (45,-1) {4};
  \node[anchor=west] at (-0.5,15.5) {\iflangja{GPU スライス（計算＋メモリ）}{GPU slices (compute + memory)}};
\end{tikzpicture}

% Width: margin (50 mm); the drawing is about 48 mm wide.
\begin{tikzpicture}[x=1mm, y=1mm, font=\scriptsize\sffamily,
  g/.style={draw, fill=ln-accent!30, minimum width=6mm, minimum height=5mm, inner sep=0pt},
  m/.style={draw, fill=ln-tip!35, minimum width=6mm, minimum height=3.5mm, inner sep=0pt},
  app/.style={draw, fill=ln-box, minimum height=4.5mm, inner sep=0pt}]
  \foreach \i in {0,...,6} {
    \node[g] at (3+\i*7,10) {\the\numexpr\i+1\relax};
    \node[m] at (3+\i*7,5) {};
  }
  \node[app, minimum width=20mm] at (10,-1) {App 1};
  \node[app, minimum width=13mm] at (27.5,-1) {App 2};
  \node[app, minimum width=6mm] at (38,-1) {3};
  \node[app, minimum width=6mm] at (45,-1) {4};
  \node[anchor=west] at (-0.5,15.5) {\iflangja{GPU スライス（計算＋メモリ）}{GPU slices (compute + memory)}};
\end{tikzpicture}

% Width: margin (50 mm); the drawing is about 48 mm wide.
\begin{tikzpicture}[x=1mm, y=1mm, font=\scriptsize\sffamily,
  g/.style={draw, fill=ln-accent!30, minimum width=6mm, minimum height=5mm, inner sep=0pt},
  m/.style={draw, fill=ln-tip!35, minimum width=6mm, minimum height=3.5mm, inner sep=0pt},
  app/.style={draw, fill=ln-box, minimum height=4.5mm, inner sep=0pt}]
  \foreach \i in {0,...,6} {
    \node[g] at (3+\i*7,10) {\the\numexpr\i+1\relax};
    \node[m] at (3+\i*7,5) {};
  }
  \node[app, minimum width=20mm] at (10,-1) {App 1};
  \node[app, minimum width=13mm] at (27.5,-1) {App 2};
  \node[app, minimum width=6mm] at (38,-1) {3};
  \node[app, minimum width=6mm] at (45,-1) {4};
  \node[anchor=west] at (-0.5,15.5) {\iflangja{GPU スライス（計算＋メモリ）}{GPU slices (compute + memory)}};
\end{tikzpicture}

  \caption{Four applications on a GPU of seven MIG slices, with GPU
    instances of 3, 2, 1 and 1 slices.  Memory simplified: the A100 has
    eight memory slices, and a 3-slice instance gets four
    \cite[§6.2]{nvidia2026mig}.}
  \label{fig:l09-mig}
\end{marginfigure}

MIG thus makes one large GPU an efficient host for several jobs, with
memory and core isolation for several users.  The A100 \enquote{can be
securely partitioned into as many as seven separate GPU Instances}
\cite{nvidia2020ampere}, which is what a cloud provider needs to rent
parts of one GPU to different customers (\cref{fig:l09-mig}).\margin{An
A100 GPU instance is built from \enquote{GPU slices} of one GPC with 14
SMs and one group of ten L2 slices; the GPU has seven
\cite{nvidia2020ampere}.  Profiles run from 1g.5gb to 7g.40gb on an
A100 \SI{40}{\giga\byte} \cite[§6.2]{nvidia2026mig}.  The H100's second-generation MIG keeps seven
instances and adds confidential computing \cite{nvidia2022hopper}.}  The
price of isolation is rigidity: an instance has a fixed size, and work
cannot spill into an idle neighbour.

\subsection{Multi-Process Service}

\begin{definition}[Time slicing, spatial partitioning, MPS]
  Under \term{time slicing}[タイムスライシング] the GPU runs the work of
  one process (one CUDA context) at a time and switches between
  processes;\caution{Lecture vs document: MPS does not give each client
  its own SMs.  A provisioning limit caps a client's share and reserves
  nothing; \enquote{kernels launched from different MPS client contexts
  may execute on the same SM}.  Only a non-uniform split (say 30\,\% and
  70\,\%), or the static SM partitions of Ampere and later, confine the
  clients to disjoint SMs \cite[§1.1.2.2.5, §1.1.2.2.11]{nvidia2026mps}.} under \term{spatial partitioning}[空間分割] the work of
  several processes runs at the same time on disjoint parts of the GPU,
  such as different SMs.  The
  \term{Multi-Process Service (MPS)}[マルチプロセスサービス（MPS）] is
  NVIDIA's runtime service that lets the CUDA work of several processes
  run concurrently on one GPU, optionally with a per-client limit on the
  SMs it may use.
  \index{MPS|see{Multi-Process Service (MPS)}}%
\end{definition}

Kernels of different CUDA contexts cannot run concurrently without MPS:
\enquote{The GPU may time slice to provide forward progress to each
context} (CPG §3.2.8.2 Concurrent Kernel Execution).  With MPS, which
supports spatial partitioning from the V100 on, the clients' kernels
share the GPU at once, but the resources are not strictly partitioned:
the L2 and the DRAM bandwidth are shared, so a client's performance
depends on what the others do and can be unpredictable.\caution{Before
Volta, MPS already let clients run concurrently, through a server that
merged their work into one context; Volta added hardware support so that
\enquote{MPS client processes have fully isolated GPU address spaces}
and \enquote{limited execution resource provisioning}, the per-client
SM share (environment variable
\texttt{CUDA\_MPS\_\allowbreak ACTIVE\_\allowbreak THREAD\_\allowbreak PERCENTAGE}) \cite{nvidia2026mps}.}
MPS also makes it easy to bring MPI programs to CUDA: the MPI ranks on
one node are separate processes, and with MPS their kernels can fill one
GPU together.  When QoS is not critical, MPS can achieve better
utilization than MIG, because a client's share is not rounded to whole
slices and idle SMs can be given to whoever has work.

\begin{example}[Fitting three jobs]
  \label{ex:l09-fit}
  Three inference jobs need 30\,\%, 25\,\% and 10\,\% of a GPU's SMs.
  Under MIG with seven equal slices, a slice is $1/7 \approx 14.3\,\%$,
  so the jobs need 3, 2 and 1 slices (42.9\,\%, 28.6\,\% and 14.3\,\%):
  six of the seven slices are taken although only 65\,\% of the SMs do
  useful work, and a fourth job that needs 20\,\% no longer fits.  Under
  MPS each job can be limited to its share
  (\texttt{CUDA\_MPS\_\allowbreak ACTIVE\_\allowbreak THREAD\_\allowbreak PERCENTAGE}
  of 30, 25 and 10, rounded down to a supported thread count).  Nothing
  is reserved, so the remaining capacity stays available to further
  clients, at the cost of sharing SMs as well as the L2 and the DRAM
  bandwidth, and no job is guaranteed its performance
  \cite[§1.1.2.2.5]{nvidia2026mps}.
\end{example}

\caution{MPS isolation is limited in a second way.  A fatal GPU fault of
one Volta MPS client is reported to all clients that share the GPUs on
which it occurred, \enquote{without indicating which client generated
the error}, and those clients must exit \cite{nvidia2026mps}.  Under MIG
a fault stays inside its instance \cite{nvidia2020ampere}.}

\subsection{Stream-based programming}

Inside one process, streams are the means of concurrency.

\begin{definition}[Concurrent kernel execution]
  The execution of several kernels of one CUDA context on the GPU at
  the same time is called
  \term{concurrent kernel execution}[カーネルの並行実行].  Kernels issued to different streams may run concurrently
  when resources allow; kernels in the same stream run in order (CPG
  §3.2.8.2 Concurrent Kernel Execution).
\end{definition}

\margin{Up to 128 kernels can be resident at once on V100, A100 and
H100 (\enquote{maximum number of resident grids per device}, CPG
Table~21 \cite{nvidia2024cudaguide}).}%
Streams thus allow concurrent kernels, the overlap of communication and
computation across streams (Lesson~5), and the submission of work by
several CPU threads, each with its own streams.  The loop below splits a
job into more chunks than there are streams and deals the chunks out to
the streams in turn; each iteration issues one chunk's copy-in, kernel
and copy-out before moving to the next chunk, a \emph{depth-first}
launch order.

\begin{lstlisting}[language=C++, numbers=none]
// h_x, h_y pinned; d_x, d_y device buffers of n floats
const int nStreams = 4, nChunks = 16;
size_t chunk = n / nChunks;
size_t bytes = chunk * sizeof(float);
cudaStream_t st[nStreams];
for (int k = 0; k < nStreams; ++k)
  cudaStreamCreate(&st[k]);
for (int i = 0; i < nChunks; ++i) {    // depth-first
  cudaStream_t s = st[i % nStreams];
  size_t off = i * chunk;
  cudaMemcpyAsync(d_x + off, h_x + off, bytes,
                  cudaMemcpyHostToDevice, s);
  square<<<(chunk + 255) / 256, 256, 0, s>>>(
      d_x + off, d_y + off, chunk);
  cudaMemcpyAsync(h_y + off, d_y + off, bytes,
                  cudaMemcpyDeviceToHost, s);
}
cudaDeviceSynchronize();               // all streams
\end{lstlisting}

Chunks 0, 4, 8 and 12 share stream~0 and therefore run one after the
other, while the four streams overlap.\tip{An OpenMP program ports
naturally to this style: a parallel loop whose iterations are
independent becomes a loop that issues each iteration's work to its own
stream, and each OpenMP thread can also drive a stream of its own.}
The scheduling has a cost: for every operation the driver prepares the
launch and the hardware must track the dependencies between operations
of the same stream.  For short kernels this overhead can be a
significant fraction of the run time.

\begin{definition}[CUDA Graph]
  A \term{CUDA Graph}[CUDAグラフ] is a set of operations (kernel
  launches, copies and others) connected by dependencies, defined
  separately from its execution: the graph is built once, instantiated
  into an executable graph that is validated and set up in advance, and
  then launched as often as needed with one call (CPG §3.2.8.7 CUDA
  Graphs).
\end{definition}

\begin{lstlisting}[language=C++, numbers=none]
cudaGraph_t g;  cudaGraphExec_t ge;
cudaStreamBeginCapture(s, cudaStreamCaptureModeGlobal);
for (int k = 0; k < 20; ++k)          // 20 short kernels
  step<<<blocks, 256, 0, s>>>(d_u, n);
cudaStreamEndCapture(s, &g);          // record, not run
cudaGraphInstantiate(&ge, g, 0);      // set up once
for (int it = 0; it < 1000; ++it)
  cudaGraphLaunch(ge, s);             // one call each
cudaStreamSynchronize(s);
\end{lstlisting}
Stream capture turns existing stream code into a graph: between
\texttt{cudaStream\allowbreak BeginCapture} and
\texttt{cudaStream\allowbreak EndCapture} the
work issued to the stream is recorded instead of executed.  Because the
whole dependency chain is known before the first launch, the setup is
paid once rather than for every kernel, and the runtime can optimise the
workflow as a whole (CPG §3.2.8.7).\caution{Stream capture does not
work on the legacy default stream (CPG §3.2.8.7.3 Creating a Graph
Using Stream Capture); \texttt{s} must not be the legacy default
stream: a created stream or \texttt{cudaStream\allowbreak PerThread}
will do.}

\subsection{Programming several GPUs}

The host program drives several GPUs by switching the current device:
the call \texttt{cudaSetDevice(d)} makes GPU~\texttt{d} current for the calling
host thread: allocations and kernel launches go to the current device,
and streams and events belong to the device that was current when they
were created (CPG §3.2.9.2 Device Selection).  The program below divides
an array among all GPUs, the data decomposition of Lesson~2, with a
separate allocation, copy and kernel launch per GPU.\margin{With a pageable \texttt{h\_x}, each
\texttt{cudaMemcpyAsync} first copies into a staging buffer on the CPU
(Lesson~5), and the GPUs start one after the other; pinned host memory
lets all copies overlap.}

\begin{lstlisting}[language=C++, numbers=none]
int nDev;
cudaGetDeviceCount(&nDev);
size_t part = n / nDev;               // n divisible
size_t bytes = part * sizeof(float);
float *d_x[16];                        // nDev <= 16
for (int d = 0; d < nDev; ++d) {
  cudaSetDevice(d);                    // select GPU d
  cudaMalloc(&d_x[d], bytes);          // on GPU d
  cudaMemcpyAsync(d_x[d], h_x + d * part, bytes,
                  cudaMemcpyHostToDevice);
  square<<<(part + 255) / 256, 256>>>(d_x[d], d_x[d],
                                      part);
}
for (int d = 0; d < nDev; ++d) {       // wait for all
  cudaSetDevice(d);
  cudaDeviceSynchronize();
}
\end{lstlisting}

\margin{A DGX H100 node holds eight H100 GPUs joined through four
NVSwitches, and eight \SI{400}{\giga\bit\per\second} ConnectX-7 network
adapters (HWP App.~A \cite{nvidia2022hopper}).}%
Kernel launches and asynchronous copies return at once, so the first
loop starts the work on every GPU before any of it has finished; the GPUs
then run in parallel, each on its own part.\caution{\texttt{cudaDeviceSynchronize}
waits only for the \emph{current} device.  One call after the launch
loop waits for the last GPU alone; the host must select each device in
turn and synchronise it, as above.}  Each device has its own default
stream (CPG §3.2.9.3 Stream and Event Behavior), and streams can be
combined with several GPUs: each GPU receives its own streams, created
while it is current.

\subsection{Comparison}

\Cref{tab:l09-mechanisms} summarises the three mechanisms.\caution{Lecture
vs document: 48 MPS clients is the Volta (CC~7.0) limit of earlier
documentation, 16 before Volta \cite{nvidia2025mpscontrol}.  The current
overview allows up to 60 per device and server with the default
\texttt{CUDA\_\allowbreak DEVICE\_\allowbreak MAX\_\allowbreak CONNECTIONS}
of 2, fewer if it is raised \cite[§1.1.2.2.6]{nvidia2026mps}.}  They are
not exclusive: a process inside one MIG instance can still use streams,
and MPS can run inside a MIG instance, with the client limit lowered in
proportion to the instance size \cite[§7.2]{nvidia2026mig}.

\begin{table}[tb]
  \centering
  \small
  \caption{Streams, MPS and MIG compared.}
  \label{tab:l09-mechanisms}
  \begin{tabular}{@{}>{\raggedright\arraybackslash}p{38mm}
                  >{\raggedright\arraybackslash}p{20mm}
                  >{\raggedright\arraybackslash}p{27mm}
                  >{\raggedright\arraybackslash}p{18mm}@{}}
    \toprule
    & Streams & MPS & MIG \\
    \midrule
    Partition type & none & logical & physical \\
    Maximum partitions & unlimited & 48 (see caution) & 7 \\
    SM isolation & no & by percentage (a cap) & yes \\
    Memory bandwidth isolation & no & no & yes \\
    Performance QoS & no & partial & yes \\
    Reconfiguration & dynamic & at process launch & when idle \\
    \bottomrule
  \end{tabular}
\end{table}

Streams cost nothing to set up and adapt at run time but give no
isolation; MIG gives the strongest isolation but changes its partition
only when the GPU is idle; MPS lies in between, with SM limits normally
set when a client process starts (per-context limits can also be set
later \cite[§1.1.2.2.5]{nvidia2026mps}).\tip{Use MIG when QoS matters,
MPS when utilization matters, and streams for concurrency inside one
process.}

\begin{keypoint}[Lesson summary]
  A monolithic GPU is limited by die size and yield, so GPUs scale as
  multi-chip modules of GPU modules, multi-GPU nodes and clusters.
  Within a package or node the memory is one NUMA space of near and far
  memory: first-touch placement, locality-aware thread block scheduling
  and page migration keep accesses near, and NVLink and NVSwitch make
  far accesses cheaper than PCIe.  Across nodes the memory is
  distributed and RDMA, with GPUDirect RDMA for GPU buffers, moves data
  without the CPUs.  To raise the utilization of one GPU, streams run
  independent work of one process without partitioning; MPS runs several
  processes at once, with optional per-client SM limits, but shares the
  L2 and DRAM; MIG splits SMs,
  L2 and DRAM into up to seven isolated instances with QoS.  CUDA Graphs
  remove the per-launch overhead of streams, and \texttt{cudaSetDevice}
  selects the GPU for each allocation, copy and launch.
\end{keypoint}

\end{document}
